\documentclass[12pt]{amsart}
%\pagestyle{plain}

\usepackage{amsmath,amssymb,amsfonts}
\usepackage{mathrsfs}
\usepackage{mathtools}
\usepackage{enumitem}
\usepackage{xcolor}
\usepackage{tikz-cd}
\setcounter{MaxMatrixCols}{30}
\setcounter{page}{1}
\usepackage{hyperref}
\usepackage[alphabetic,msc-links]{amsrefs}
\hypersetup{
  colorlinks=false,
  pdfborder={0 0 0.7},
  pdfborderstyle={/S/D/D[3 2]/W 0.7},
  linkbordercolor={0.10 0.52 0.28},
  citebordercolor={0.04 0.45 0.22},
  urlbordercolor={0.00 0.42 0.30},
  pdftitle={McKean Rigidity for Cocompact Negatively Curved Manifolds and the p-Laplacian},
  pdfauthor={Kuntao Jin, Shing-Tung Yau, and Bo Zhu},
  pdfsubject={Spectral rigidity under an upper sectional-curvature bound},
  pdfkeywords={McKean inequality, p-Laplacian, bottom spectrum, Busemann function, leafwise diffusion}
}

% The msc-links option makes every MR number a MathSciNet hyperlink.

\usepackage{geometry}
\geometry{left=3cm,right=3cm,top=3cm,bottom=3.5cm}
\setcounter{section}{0}

% Theorem environments
\newtheorem{theorem}{Theorem}[section]
\newtheorem{conjecture}[theorem]{Conjecture}
\newtheorem{proposition}[theorem]{Proposition}
\newtheorem{lemma}[theorem]{Lemma}
\newtheorem{corollary}[theorem]{Corollary}
\newtheorem{question}[theorem]{Question}
\newtheorem{problem}[theorem]{Problem}
\newtheorem{fact}[theorem]{Fact}
\newtheorem{claim}[]{Claim}
\newtheorem*{acknowledgements}{Acknowledgements}
\newtheorem*{quotedtheorem}{Theorem}
\newtheorem*{hilleyosida}{Hille--Yosida theorem (contraction form)}
\theoremstyle{definition}
\newtheorem{definition}[theorem]{Definition}
\newtheorem{example}[theorem]{Example}
\newtheorem{remark}[theorem]{Remark}
\newtheorem*{unnumberedremark}{Remark}

\numberwithin{equation}{section}
\newcommand{\dv}{\mathrm{div}}
\newcommand{\ee}{{\mathrm{e}}}
\newcommand{\lb}[1]{\langle#1\rangle}
\newcommand{\Blb}[1]{\Big\langle#1\Big\rangle}
\newcommand{\blb}[1]{\big\langle#1\big\rangle}
\newcommand{\mf}{\mathbf}
\newcommand{\mb}{\mathbb}
\newcommand{\mc}{\mathcal}
\newcommand{\ms}{\mathscr}
\newcommand{\mk}{\mathfrak}
\newcommand{\oli}{\overline}
\newcommand{\eps}{\varepsilon}
\newcommand{\wti}{\widetilde}
\newcommand{\mr}{\mathrm}
\newcommand{\dist}{\operatorname{dist}}
\newcommand{\R}{\mathbb{R}}
\newcommand{\bH}{\mathbb{H}}
\newcommand{\supp}{\mathrm{supp}}
\newcommand{\vol}{\operatorname{vol}}
\newcommand{\area}{\operatorname{area}}
\newcommand{\tr}{\mathrm{tr}}
\DeclareMathOperator{\Ric}{Ric}
\DeclareMathOperator{\Hess}{Hess}

\begin{document}

\title[McKean rigidity]{McKean Rigidity for Cocompact Negatively Curved
Manifolds and the \(p\)-Laplacian}
\date{\today}

\author{Kuntao Jin}
\address[Kuntao Jin]{Department of Mathematical Sciences, Tsinghua University}
\email{jkt25@mails.tsinghua.edu.cn}


\author{Bo Zhu}
\address[Bo Zhu]{Yau Mathematical Sciences Center,
Tsinghua University}
\email{zhub@tsinghua.edu.cn}
\thanks{Bo Zhu is supported by NSFC~12501066.}



\subjclass[2020]{Primary 53C24, 58J50; Secondary 35J92, 58J65}
\keywords{\(p\)-Laplacian, McKean inequality, bottom spectrum, Busemann function,
horospherical suspension, leafwise diffusion, spectral rigidity}



\begin{abstract}
Let \((M^m,g)\) be a closed Riemannian manifold with
\(\sec_g\leq-1\).  We prove that the bottom spectrum of its universal
cover attains McKean's lower bound if and only if the universal cover is
hyperbolic space of constant sectional curvature \(-1\).  More generally,
for every \(1<p<\infty\), the variational \(p\)-fundamental tone satisfies
\[
  \lambda_{1,p}(\wti M)
  \geq\left(\frac{m-1}{p}\right)^p,
\]
and equality for some \(p\in(1,\infty)\) holds if and only if
\(\wti M\cong\bH^m(-1)\).  In that case, equality holds for every
\(p\in(1,\infty)\).  The proof converts the two McKean defects of a
minimizing sequence into a stationary probability measure on the compact
horospherical suspension; heat-kernel positivity then forces its
zero-defect support to contain a complete leaf. 
\end{abstract}
\maketitle
%\tableofcontents
\section{Introduction}
Let \((X^m,g_X)\) be a complete Riemannian manifold, where \(m\geq2\).  Its
bottom spectrum is
\begin{equation}\label{eq:bottom-spectrum}
  \lambda_1(X)
  =\inf_{0\neq u\in C_c^\infty(X)}
  \frac{\int_X|\nabla u|^2\,dV}
       {\int_Xu^2\,dV}.
\end{equation}
If \(X\) is simply connected and \(\sec_{g_X}\leq-1\), McKean proved the
following sharp estimate (see \cite{McKean1970}, pp.~360--365):
\begin{equation}\label{eq:mckean-bound}
  \lambda_1(X)
  \geq\frac{(m-1)^2}{4}
  =\lambda_1\bigl(\bH^m(-1)\bigr).
\end{equation}
For a general complete, simply connected manifold, however, equality in
McKean's estimate is not rigid: the bottom spectrum alone does
not determine the geometry of \(X\).  The examples in
Subsection~\ref{subsec:role-cocompactness} explain the two points at which
cocompactness enters the rigidity argument.

\medskip
We prove rigidity in the cocompact setting: \(X=\wti M\) is the universal
Riemannian cover of a closed manifold.

\begin{theorem}\label{thm:main}
Let \((M^m,g)\) be a closed connected Riemannian manifold of dimension
\(m\geq2\) such that
\(
  \sec_g\leq-1.
\)
Then
\[
  \lambda_1(\wti M)\geq\frac{(m-1)^2}{4}.
\]
Equality holds if and only if
\[
  (\wti M,\wti g)\cong\bH^m(-1).
\]
\end{theorem}

\medskip
The same rigidity holds for the variational \(p\)-Laplacian.  For a
complete Riemannian manifold \(X\) and \(1<p<\infty\), set
\[
  \lambda_{1,p}(X)
  :=
  \inf_{0\neq u\in C_c^\infty(X)}
  \frac{\|\nabla u\|_{L^p(X)}^p}
       {\|u\|_{L^p(X)}^p}.
\]
The sharp lower bound under \(\sec\leq-1\) is known
(see \cite{Poliquin2014}, Proposition~4.5, and
\cite{CarvalhoCavalcante2022}, Corollary~1.1).  Our concern is rigidity
in the equality case for universal covers of closed manifolds.

\begin{theorem}
\label{thm:p-rigidity}
Let \((M^m,g)\) be a closed connected Riemannian manifold of dimension
\(m\geq2\) satisfying \(\sec_g\leq-1\).  Then, for every
\(1<p<\infty\),
\[
  \lambda_{1,p}(\wti M)
  \geq\left(\frac{m-1}{p}\right)^p.
\]
Equality holds for some \(p\in(1,\infty)\) if and only if
\[
  (\wti M,\wti g)\cong\bH^m(-1).
\]
In this case, equality holds for every \(p\in(1,\infty)\).
\end{theorem}

To the best of our knowledge, neither the quadratic equality-rigidity
statement under the upper curvature bound \(\sec_g\leq-1\) nor its
\(p\)-Laplacian extension was previously known.  In this paper, we first
treat the quadratic case \(p=2\), where the square identity separates the
two McKean defects.  For general \(p\), uniform convexity of
\(Y\mapsto|Y|^p\) yields, for a normalized \(p\)-minimizing sequence
\(u_j\) and \(\rho_j=u_j^p\),
\[
  \left\|
    \nabla\rho_j+(m-1)\rho_j\nabla B
  \right\|_{L^1}
  \longrightarrow0.
\]
The coefficient \(m-1\) is independent of \(p\); hence the limiting measure
is stationary for the same leafwise diffusion as in the quadratic case,
and the geometric and dynamical parts of the argument are unchanged.  The
nonlinear defect calculation and the proof of
Theorem~\ref{thm:p-rigidity} are deferred to
Appendix~\ref{app:p-rigidity}.

\bigskip
\subsection*{Related rigidity results under curvature bounds}

For a complete Riemannian manifold \((X^m,g)\) with
\(\Ric_g\geq-(m-1)g\), Shiu-Yuen Cheng proved the sharp estimate
(see \cite{Cheng1975}):
\(\lambda_1(X)\leq (m-1)^2/4\). When $m\geq 3$,
Peter Li and Jiaping Wang analyzed the equality
case on general complete manifolds by harmonic-function methods; their
splitting theorems identify explicit warped products as the only equality
cases with more than one end (see \cite{LiWang2001} and
\cite{LiWang2002}).  Thus equality alone does not imply hyperbolicity in
the noncocompact setting.  For universal covers of closed manifolds,
however, Xiaodong Wang proved full rigidity using Martin-boundary harmonic
functions and Kaimanovich entropy: if \((M^m,g)\) is closed,
\(\Ric_g\geq-(m-1)g\), and
\(\lambda_1(\wti M)=(m-1)^2/4\), then
\(\wti M\cong\bH^m(-1)\)
(see \cite{Wang2008}, Theorem~4).

Scalar-curvature analogues follow the same pattern.  In dimension three,
Munteanu and Jiaping Wang used harmonic-function theory to prove the sharp
estimate, normalized to the hyperbolic curvature scale \(-1\):
\(\operatorname{Sc}_X\geq-6\) implies \(\lambda_1(X)\leq 1\)
for a complete \(X^3\) with finitely many ends and finite first Betti
number (see \cite{MunteanuWang2024}, Theorem~1.1, and
\cite{MunteanuWang2026}, Theorem~1.1).\footnote{Munteanu and Wang assume a
Ricci lower bound in the first paper and remove this assumption in the
second (see \cite{MunteanuWang2024} and \cite{MunteanuWang2026}).} In
higher dimensions, Jinmin Wang and Bo Zhu proved the corresponding
sharp estimate, together with rigidity in the equality case, for universal
covers of closed manifolds.  If
\(M\) is rationally essential, \(\wti M\) is spin, and \(\pi_1(M)\)
satisfies the Strong Novikov Conjecture, then
\(\operatorname{Sc}_g\geq-m(m-1)\) implies
\(\lambda_1(\wti M)\leq(m-1)^2/4\), and equality forces
\((M,g)\) to be hyperbolic
(see \cite{WangZhu2024}, Theorem~1.2, and
\cite{WangZhu2026}, Theorem~1.2).\footnote{Wang and Zhu prove the sharp
estimate in the first paper and rigidity in the second
(see \cite{WangZhu2024} and \cite{WangZhu2026}).  Their rigidity argument
uses almost-harmonic spinors to
recover the Einstein identity
\(\Ric_g=-(m-1)g\), after which Xiaodong Wang's rigidity theorem applies.}

Related rigidity results begin with stronger horospherical assumptions.
Besson, Courtois, and Hersonsky obtain hyperbolicity from geometric
conditions on horospheres, while Zimmer proves rigidity under asymptotic
harmonicity, which requires every Busemann function to have the same
constant Laplacian (see
\cites{BessonCourtoisHersonsky2025Flat,
BessonCourtoisHersonsky2025Horospherical,Zimmer2012}).  By contrast,
spectral equality here initially yields only averaged defect limits.  The
stationary-measure argument produces a single
\(\xi_*\in\partial_\infty\wti M\) satisfying
\(\Delta B_{\xi_*}=m-1\), and under the present curvature bound this
identity suffices to prove hyperbolicity.

\subsection*{The invariant perspective: isoperimetry and entropy}

The preceding results formulate rigidity through curvature inequalities.
A complementary viewpoint uses two global geometric invariants: the
Cheeger constant and the volume entropy.
The relation between isoperimetric constants and the first eigenvalue was
further developed by Yau in the compact setting
(see \cite{Yau1975Isoperimetric}).

For a complete \(X\), and for a universal cover \(\wti M\) of a closed
manifold, set
\[
\begin{aligned}
  h_{\mathrm{iso}}(X)
  &:=
  \inf_{\Omega\Subset X}
  \frac{\area(\partial\Omega)}{\vol(\Omega)},\\
  h_{\mathrm{vol}}(\wti M)
  &:=
  \lim_{R\to\infty}\frac{1}{R}\log\vol B(o,R).
\end{aligned}
\]
Then Cheeger's inequality and Brooks's volume-growth
estimate give (see \cite{Cheeger1970} and \cite{Brooks1981})
\begin{equation}\label{eq:cheeger-spectrum-entropy-chain}
  h_{\mathrm{iso}}(\wti M)^2
  \leq 4\lambda_1(\wti M)
  \leq h_{\mathrm{vol}}(\wti M)^2.
\end{equation}

Under a Ricci lower bound \(\Ric_g\geq-(m-1)g\), Bishop--Gromov comparison
gives \(h_{\mathrm{vol}}(\wti M)\leq m-1\).  Hence equality in the sharp
spectral bound forces \(h_{\mathrm{vol}}(\wti M)=m-1\), and the entropy
rigidity theorem of Ledrappier--X.~Wang, with Gang Liu's short proof, yields
\(\wti M\cong\bH^m(-1)\)
(see \cite{LedrappierWang2010}, Theorem~2, and
\cite{Liu2011}, Theorem~1).

Under the sectional-curvature assumption of this paper, G\"unther's volume
comparison theorem gives the opposite inequality (see \cite{Gunther1960}):
\[
  \sec_g\leq-1
  \quad\Longrightarrow\quad
  h_{\mathrm{vol}}(\wti M)\geq m-1.
\]
Together with the right-hand inequality in
\eqref{eq:cheeger-spectrum-entropy-chain}, spectral equality again gives
only \(h_{\mathrm{vol}}(\wti M)\geq m-1\), which is precisely the
comparison estimate above.  Thus the spectral hypothesis supplies no
reverse volume-growth inequality, and the entropy route does not close
the argument.

Cheeger's inequality yields a necessary consequence of spectral equality,
but it does not by itself imply rigidity.  Indeed, for every relatively
compact smooth domain \(\Omega\subset\wti M\), Busemann comparison and the
divergence theorem give
\[
  (m-1)\vol(\Omega)
  \leq\int_\Omega\Delta B_\xi\,dV
  =\int_{\partial\Omega}\langle\nabla B_\xi,\nu\rangle\,dA
  \leq\area(\partial\Omega).
\]
Taking the infimum over \(\Omega\) yields
\(h_{\mathrm{iso}}(\wti M)\geq m-1\).  If spectral equality holds, Cheeger's
inequality gives \(h_{\mathrm{iso}}(\wti M)\leq m-1\), and therefore
\[
  h_{\mathrm{iso}}(\wti M)=m-1.
\]
For \(p=1\), define
\[
  \lambda_{1,1}(X)
  :=
  \inf_{0\neq u\in C_c^\infty(X)}
  \frac{\int_X|\nabla u|\,dV}
       {\int_X|u|\,dV}.
\]
The coarea formula and smooth approximation of characteristic functions
give
\[
  \lambda_{1,1}(X)=h_{\mathrm{iso}}(X).
\]


The proof of Theorem~\ref{thm:main} uses spectral minimizing sequences
directly.  Entropy equality is not assumed; it follows only after
hyperbolicity has been proved.

\subsection*{Outline of the proof}

Assume that \(\lambda_1(\wti M)=(m-1)^2/4\).  To prove
Theorem~\ref{thm:main}, it is enough to find
\(\xi_*\in\partial_\infty\wti M\) such that \(B=B_{\xi_*}\) satisfies
\begin{equation}\label{eq:busemann-conclusion}
  |\nabla B|=1,
  \quad
  \Delta B=m-1
\end{equation}
throughout \(\wti M\).  The unit-gradient identity is automatic for a
Busemann function; the task is therefore to obtain the constant-Laplacian
identity.  Busemann Hessian comparison then converts
\eqref{eq:busemann-conclusion} into hyperbolic rigidity.

The central point is that the two McKean defects control complementary
properties of a single limiting measure on the compact horospherical
suspension.  The first defect gives stationarity for a leafwise diffusion
with drift coefficient \(m-1\), while the second confines the support to
the zero set of the Busemann Laplacian defect.  Positivity of the leafwise
diffusion makes the stationary support leaf-saturated, thereby propagating
the Busemann equality along a complete leaf.  This stationary-measure
viewpoint is inspired by Ledrappier--X.~Wang
(see \cite{LedrappierWang2010}, Theorems~1 and~2).  The required diffusion
theory follows Candel's Hille--Yosida construction
(see \cite{Candel2003}, Theorem~4.14 and
Propositions~4.15--4.16); its detailed verification is deferred to
Section~\ref{sec:candel-semigroup-construction}, after the main rigidity
argument.

\begin{enumerate}[label=\textup{(\arabic*)}]
\item Section~\ref{sec:busemann-defects} derives the McKean square
identity and produces a normalized minimizing sequence \(u_j\) for which
both defects vanish; see \eqref{eq:first-defect} and
\eqref{eq:second-defect}.

\item Section~\ref{sec:suspension} introduces the compact horospherical
suspension
\[
  Z=(\wti M\times\partial_\infty\wti M)/\Gamma,
  \quad Z\cong SM.
\]
Its leaves are covered by \(\wti M\), and its compactness prevents the
quadratic mass of \(u_j\) from being lost at infinity.

\item In Section~\ref{sec:limiting-measure}, we push \(u_j^2dV\) to \(Z\) along
the leaf determined by \(\xi_0\).  A subsequence converges to a probability
measure \(\mu\).  For
\(D([x,\xi])=\Delta B_\xi(x)-(m-1)\), the second defect gives
\[
  \supp\mu\subseteq D^{-1}(0).
\]

\item In Section~\ref{sec:stationarity}, the first defect gives
\(\int_ZLf\,d\mu=0\) for the leafwise drift operator
\[
  L=\Delta^{\mathrm{leaf}}
    -(m-1)\langle V,\nabla^{\mathrm{leaf}}\,\cdot\,\rangle.
\]
The associated diffusion makes \(\mu\) stationary, and its leafwise
positivity makes \(\supp\mu\) leaf-saturated.  Since
\(\supp\mu\subseteq D^{-1}(0)\), one complete leaf satisfies \(D=0\);
hence \(\Delta B_{\xi_*}=m-1\) on \(\wti M\).

\item In Section~\ref{sec:rigidity}, equality in Busemann Hessian comparison
gives
\[
  \nabla^2B=\wti g-dB\otimes dB,
\]
so the gradient flow yields an exponential warped product.  The curvature
bound inherited from the compact quotient forces its cross-section to be
flat, and therefore \((\wti M,\wti g)\cong\bH^m(-1)\).

\item After the main proof, Section~\ref{sec:candel-semigroup-construction}
develops the leafwise semigroup used in
Section~\ref{sec:stationarity}.  It gives the construction at the finite
transverse regularity needed here, proves gradient control on the
classical domain and the core property of that domain, and establishes
leaf saturation of stationary supports.
\end{enumerate}

Appendix~\ref{app:p-rigidity} proves
Theorem~\ref{thm:p-rigidity}.  The uniform-convexity defect for the
\(p\)-energy yields an \(L^1\) transport equation for the probability
density \(u_j^p\).  Since its drift coefficient is again \(m-1\), the
suspension, stationary-measure, and support-saturation arguments from the
quadratic case apply without change.

\begin{acknowledgements}
The authors are grateful to Jiaping Wang for bringing this problem to their attention and to Shing-Tung Yau and Xiaodong Wang for their valuable comments and suggestions. The first author also thanks his advisor, Guoyi Xu, for his guidance and mentorship.
\end{acknowledgements}

\section{Busemann functions and the McKean defects}
\label{sec:busemann-defects}

No cocompact group action is used in this section.  Throughout, let
\((X^m,g_X)\) be a complete simply connected Riemannian manifold, where
\(m\geq2\), and assume that, for some \(A\geq1\),
\begin{equation}\label{eq:pinched-curvature}
  -A^2\leq\sec_{g_X}\leq-1.
\end{equation}
Then \(X\) is a pinched negatively curved Hadamard manifold.  The
facts about the visual boundary and Busemann functions used in this
section are classical
(see \cite{EberleinONeill1973}, Sections~1--3, and
\cite{HeintzeImHof1977}, Sections~2--3).  The excess in McKean's
inequality will be written as the sum of a first-order defect and a
Busemann Laplacian defect.  Spectral equality forces both terms to vanish
along a single normalized minimizing sequence.

\subsection{The visual boundary and Busemann comparison}

A \emph{geodesic ray} is a unit-speed geodesic
\(c\colon[0,\infty)\to X\); thus
\[
  d\bigl(c(s),c(t)\bigr)=|s-t|
  \quad (s,t\geq0).
\]
Two geodesic rays \(c_1\) and \(c_2\) are called \emph{asymptotic}, written
\(c_1\sim c_2\), if
\[
  \sup_{t\geq0}d\bigl(c_1(t),c_2(t)\bigr)<\infty.
\]
For unit-speed rays, this is equivalent to requiring their images to have
finite Hausdorff distance.  The triangle inequality shows that \(\sim\) is
an equivalence relation.  The \emph{visual boundary} of \(X\) is
\[
  \partial_\infty X
  :=
  \{\text{geodesic rays in \(X\)}\}/\mathord{\sim}.
\]
The class \([c]\), also denoted by \(c(+\infty)\), is the endpoint of \(c\)
at infinity.

We next describe this boundary from a fixed base point \(o\in X\).  Set
\[
  S_oX:=\{v\in T_oX:|v|_{g_X}=1\},
  \quad
  c_v(t):=\exp_o(tv).
\]
Every \(\xi\in\partial_\infty X\) has a unique representative
\(c_{o,\xi}\) issuing from \(o\).  Existence follows by taking a locally
uniform limit of the minimizing segments from \(o\) to points tending to
\(\xi\).  For uniqueness, if \(\gamma_1,\gamma_2\) are asymptotic rays
issuing from \(o\), then
\[
  f(t):=d\bigl(\gamma_1(t),\gamma_2(t)\bigr)
\]
is bounded and convex with \(f(0)=0\).  Thus
\(f(s)\leq (s/t)f(t)\) for \(0<s<t\); letting \(t\to\infty\) gives
\(\gamma_1=\gamma_2\).

It follows that the endpoint map
\[
  \Phi_o\colon S_oX\longrightarrow\partial_\infty X,
  \quad
  \Phi_o(v):=[c_v],
\]
is a bijection.  The \emph{cone topology} on \(\partial_\infty X\) is the
topology transported from \(S_oX\) by \(\Phi_o\).  Thus
\(\xi_j\to\xi\) precisely when
\[
  \dot c_{o,\xi_j}(0)\longrightarrow\dot c_{o,\xi}(0)
  \quad\text{in }S_oX.
\]
Equivalently, \(c_{o,\xi_j}\to c_{o,\xi}\) uniformly on every bounded time
interval.

For another base point \(o'\in X\), define the reanchoring map by
\[
  R_{o,o'}
  :=
  \Phi_{o'}^{-1}\circ\Phi_o
  \colon S_oX\longrightarrow S_{o'}X.
\]
The same compactness and convexity argument shows that \(R_{o,o'}\) and
\(R_{o',o}\) are continuous.  Hence \(R_{o,o'}\) is a homeomorphism, so the
cone topology is independent of the base point.

Finally, an isometry \(F\in\operatorname{Isom}(X)\) acts on the visual
boundary by
\[
  F_\infty([c]):=[F\circ c].
\]
This is well defined because \(F\) preserves distances.  In the endpoint
coordinates it has the form
\[
  F_\infty
  =
  \Phi_{F(o)}
  \circ\left.dF_o\right|_{S_oX}
  \circ\Phi_o^{-1},
\]
so \(F_\infty\) is a homeomorphism of \(\partial_\infty X\).

For \(\xi\in\partial_\infty X\), define the Busemann function normalized at
\(o\) by
\begin{equation}\label{eq:busemann-definition}
  B_{\xi,o}(x)
  :=\lim_{t\to\infty}
  \bigl(d(x,c_{o,\xi}(t))-t\bigr).
\end{equation}
To justify the definition, set
\[
  h_t(x):=d(x,c_{o,\xi}(t))-t.
\]
The triangle inequality shows that \(t\mapsto h_t(x)\) is nonincreasing and
that
\[
  -d(x,o)\leq h_t(x)\leq d(x,o).
\]
Hence the limit in \eqref{eq:busemann-definition} exists, is finite, and
satisfies \(B_{\xi,o}(o)=0\).  Since the functions \(h_t\) are uniformly
\(1\)-Lipschitz, they converge uniformly on compact subsets, and
\(B_{\xi,o}\) is \(1\)-Lipschitz.

Changing the base point changes only the additive normalization.  More
precisely, for \(o'\in X\),
\[
  B_{\xi,o'}(x)
  =
  B_{\xi,o}(x)-B_{\xi,o}(o').
\]
It is therefore natural to introduce the Busemann cocycle
\[
  \beta_\xi(x,y)
  :=
  B_{\xi,o}(x)-B_{\xi,o}(y);
\]
by the base-point change identity, this definition is independent of \(o\).
It satisfies
\[
  \beta_\xi(x,z)
  =
  \beta_\xi(x,y)+\beta_\xi(y,z),
  \quad
  |\beta_\xi(x,y)|\leq d(x,y).
\]
Moreover, if \(F\in\operatorname{Isom}(X)\), then
\[
  B_{F_\infty(\xi),F(o)}(F(x))=B_{\xi,o}(x),
  \quad
  \beta_{F_\infty(\xi)}(F(x),F(y))=\beta_\xi(x,y).
\]
Since the gradient, Hessian, and Laplacian are unchanged by adding a
constant, we suppress the base point and write \(B_\xi\) whenever only
derivatives are involved.  The curvature pinching in
\eqref{eq:pinched-curvature} yields the differential comparison estimates
used below.

\begin{proposition}
\label{prop:busemann-comparison}
Let \((X^m,g_X)\) be a complete simply connected Riemannian manifold,
where \(m\geq2\), and assume that
\[
  -A^2\leq\sec_{g_X}\leq-1
\]
for some \(A\geq1\).
Then 
\begin{enumerate}[label=\textup{(\roman*)}]
\item For every \(\xi\in\partial_\infty X\), the Busemann function
  \(B_\xi\) is of class \(C^2\).  If \(v_{x,\xi}\in S_xX\) is the initial
  velocity of the unique ray from \(x\) to \(\xi\), then
  \begin{equation}\label{eq:unit-busemann-gradient}
    \nabla B_\xi(x)=-v_{x,\xi},
    \quad
    |\nabla B_\xi(x)|=1.
  \end{equation}

\item For every \(\xi\in\partial_\infty X\),
  \addtocounter{equation}{1}
  \begin{equation}\label{eq:busemann-laplacian-comparison}
    \begin{aligned}
      g_X-dB_\xi\otimes dB_\xi
      &\leq\nabla^2B_\xi
      \leq A(g_X-dB_\xi\otimes dB_\xi),\\
      m-1&\leq\Delta B_\xi\leq A(m-1).
    \end{aligned}
  \end{equation}

\item Let \(\operatorname{pr}_X\colon
  X\times\partial_\infty X\to X\) be the projection.  Then
  \[
    (x,\xi)\longmapsto\nabla B_\xi(x),
    \quad
    (x,\xi)\longmapsto\nabla^2B_\xi(x)
  \]
  define continuous sections of
  \(\operatorname{pr}_X^*TX\) and
  \(\operatorname{pr}_X^*\operatorname{Sym}^2(T^*X)\), respectively, and
  \[
    (x,\xi)\longmapsto\Delta B_\xi(x)
  \]
  is a continuous function on \(X\times\partial_\infty X\).
\end{enumerate}
\end{proposition}

\begin{proof}
Fix \(\xi\in\partial_\infty X\) and \(x\in X\), and write
\(c=c_{x,\xi}\).  For \(R>0\), set
\[
  b_R(y):=d(y,c(R))-R.
\]
On every compact subset of \(X\), the functions \(b_R\) are smooth for
all sufficiently large \(R\) and converge to the Busemann function
normalized by \(B_\xi(x)=0\).

We use the standard radial-field convergence theorem.  The radial fields
\(-\nabla b_R\) converge locally uniformly to
\(v_{\,\cdot,\xi}\), and their covariant derivatives converge locally
uniformly to the derivatives determined by the stable Jacobi fields along
the rays to \(\xi\).  Equivalently, \(b_R\to B_\xi\) in
\(C^2\) on compact subsets.  This gives
\[
  B_\xi\in C^2(X),
  \quad
  \nabla B_\xi(x)=-v_{x,\xi},
  \quad
  |\nabla B_\xi|=1,
\]
and proves \textup{(i)}.  The local \(C^2\)-convergence, including its
uniformity on compact families, is the content of the stable-Jacobi-field
argument cited after the proof.

For \textup{(ii)}, decompose
\[
  w=a\,v_{x,\xi}+w^\perp,
  \quad
  w^\perp\perp v_{x,\xi}.
\]
Hessian comparison for the distance from \(c(R)\) gives
\[
  \coth R\,|w^\perp|^2
  \leq \nabla^2b_R(x)(w,w)
  \leq A\coth(AR)\,|w^\perp|^2.
\]
Passing to the \(C^2\)-limit and using
\(\nabla B_\xi=-v_{x,\xi}\), we obtain
\[
  |w|^2-\langle w,\nabla B_\xi\rangle^2
  \leq\nabla^2B_\xi(w,w)
  \leq
  A\bigl(|w|^2-\langle w,\nabla B_\xi\rangle^2\bigr).
\]
This is the tensor inequality in
\eqref{eq:busemann-laplacian-comparison}.  Taking the trace in an
orthonormal basis whose first vector is \(v_{x,\xi}\) gives the two
Laplacian bounds.

It remains to prove the joint continuity in \textup{(iii)}.  The map
\[
  (x,\xi)\longmapsto v_{x,\xi}
\]
is continuous by the cone topology, so the gradient formula gives the
joint continuity of \(\nabla B_\xi(x)\).

For the Hessian, fix \(T>0\) and define
\[
  H^T_{x,\xi}
  :=
  \left.
  \nabla_y^2\bigl(d(y,c_{x,\xi}(T))-T\bigr)
  \right|_{y=x}.
\]
Because a Hadamard manifold has no cut locus and
\((x,\xi)\mapsto c_{x,\xi}(T)\) is continuous, \(H^T\) is a continuous
section of
\(\operatorname{pr}_X^*\operatorname{Sym}^2(T^*X)\).

Let \(J_w^T\) be the normal Jacobi field along
\(c_{x,\xi}|_{[0,T]}\) with \(J_w^T(0)=w\) and \(J_w^T(T)=0\), and let
\(J_w^s\) be the bounded normal Jacobi field with \(J_w^s(0)=w\).
The Euclidean comparison estimate in the construction of \(J_w^s\)
gives
\[
  \bigl|(J_w^T)'(0)-(J_w^s)'(0)\bigr|
  \leq\frac{|w|}{T}.
\]
The Jacobi-field formulas for the two Hessians, together with their
vanishing in the \(v_{x,\xi}\)-direction, therefore yield
\[
  \bigl\|H^T_{x,\xi}-\nabla^2B_\xi(x)\bigr\|
  \leq\frac1T.
\]
The estimate is independent of \((x,\xi)\).  Hence
\(\nabla^2B_\xi(x)\) is the uniform limit of the continuous sections
\(H^T\) and is jointly continuous.  Taking the trace proves the joint
continuity of \(\Delta B_\xi(x)\).
\end{proof}

Proposition~\ref{prop:busemann-comparison} establishes the regularity and
comparison facts needed below.  The \(C^2\)-regularity in part~\textup{(i)}
follows from the stable-Jacobi-field argument (see
\cite{HeintzeImHof1977}, Proposition~3.1 and Lemma~2.2).  The
joint-continuity statement in part~\textup{(iii)} is also standard
(see \cite{EberleinONeill1973}, Proposition~3.5, and
\cite{HeintzeImHof1977}, Proposition~3.2).  The two-sided curvature bound
supplies the regularity used here; Busemann functions on an arbitrary
Hadamard manifold need not be \(C^2\).  In the model case of constant
sectional curvature \(-1\), the comparison estimate is sharp:
\[
  \nabla^2B_\xi=g_X-dB_\xi\otimes dB_\xi,
  \quad
  \Delta B_\xi=m-1.
\]
Thus, under \eqref{eq:pinched-curvature}, both \(\nabla^2B_\xi\) and
\(\Delta B_\xi\) are classical pointwise quantities, and the square
identity below follows by ordinary integration by parts.

\subsection{The McKean square identity}
The following identity is essential for the arguments that follow, and it is independent of curvature.
\begin{lemma}
\label{lem:mckean-square}
Let \((X^m,g_X)\) be a Riemannian manifold without boundary, where
\(m\geq2\), and let \(B\in C^2(X)\) satisfy \(|\nabla B|=1\).  Then, for every
\(u\in C_c^\infty(X)\),
\begin{equation}\label{eq:mckean-square}
\begin{split}
  \int_X\left(|\nabla u|^2-\frac{(m-1)^2}{4}u^2\right)\,dV
  ={}&\int_X\left|\nabla u+\frac{m-1}{2}u\nabla B\right|^2\,dV \\
  &+\frac{m-1}{2}\int_X(\Delta B-(m-1))u^2\,dV.
\end{split}
\end{equation}
\end{lemma}

\begin{proof}
Set
\(
  I=\int_X|\nabla u+\frac{m-1}{2}u\nabla B|^2\,dV.
\)
Since \(u\) has compact support, integration by parts has no boundary
term.  Using \(|\nabla B|=1\) and
\(\nabla(u^2)=2u\nabla u\), we obtain
\begin{align*}
  I
  ={}&\int_X|\nabla u|^2\,dV
  +\frac{m-1}{2}\int_X
     \langle\nabla(u^2),\nabla B\rangle\,dV
  +\frac{(m-1)^2}{4}\int_Xu^2\,dV \\
  ={}&\int_X|\nabla u|^2\,dV
  -\frac{m-1}{2}\int_X(\Delta B)u^2\,dV
  +\frac{(m-1)^2}{4}\int_Xu^2\,dV \\
  ={}&\int_X\left(
    |\nabla u|^2-\frac{(m-1)^2}{4}u^2
  \right)\,dV \\
  &-\frac{m-1}{2}
  \int_X\bigl(\Delta B-(m-1)\bigr)u^2\,dV.
\end{align*}
Adding the final integral to both sides gives
\eqref{eq:mckean-square}.
\end{proof}

Lemma~\ref{lem:mckean-square} is the key analytic identity of the paper.
Apply it with \(B=B_\xi\).  Proposition~\ref{prop:busemann-comparison}
gives \(|\nabla B_\xi|=1\) and
\(\Delta B_\xi-(m-1)\geq0\).  Thus
\eqref{eq:mckean-square} decomposes the excess in McKean's inequality into
two nonnegative terms and therefore directly implies
\eqref{eq:mckean-bound}.  Moreover, if equality holds in McKean's
inequality, then, for every fixed \(\xi\), any \(L^2\)-normalized
minimizing sequence makes both terms tend to zero simultaneously.  These
are precisely the first and second defect limits
\eqref{eq:first-defect} and \eqref{eq:second-defect}.

We call the squared \(L^2\)-term the \emph{first defect}; it measures the
first-order error in
\[
  \nabla u+\frac{m-1}{2}u\nabla B_\xi=0.
\]
We call the integral containing
\(\Delta B_\xi-(m-1)\) the \emph{second defect}; it is the
\(u^2\,dV\)-weighted error in the Busemann Laplacian equality
\[
  \Delta B_\xi=m-1.
\]

\subsection{Vanishing of the first and second defects}

The next proposition makes the preceding consequence of
Lemma~\ref{lem:mckean-square} explicit.

\begin{proposition}
\label{prop:defects}
Let \((X^m,g_X)\) be a complete simply connected Riemannian manifold with
\(m\geq2\) and \(-A^2\leq\sec_{g_X}\leq-1\) for some \(A\geq1\).  Assume
\[
  \lambda_1(X)=\frac{(m-1)^2}{4}.
\]
Then there exists an \(L^2\)-normalized minimizing sequence
\(\{u_j\}\subset C_c^\infty(X)\) such that
\begin{equation}\label{eq:minimizing-sequence}
  \int_Xu_j^2\,dV=1,
  \quad
  \int_X|\nabla u_j|^2\,dV
  \longrightarrow\frac{(m-1)^2}{4}.
\end{equation}
For every fixed \(\xi_0\in\partial_\infty X\), the same sequence satisfies
the following two conclusions.
\begin{enumerate}[label=\textup{(\roman*)}]
\item The first defect vanishes:
\begin{equation}\label{eq:first-defect}
  \int_X\left|
    \nabla u_j+\frac{m-1}{2}u_j\nabla B_{\xi_0}
  \right|^2\,dV\longrightarrow0.
\end{equation}
\item The second defect vanishes:
\begin{equation}\label{eq:second-defect}
  \int_X\bigl(\Delta B_{\xi_0}-(m-1)\bigr)u_j^2\,dV
  \longrightarrow0.
\end{equation}
\end{enumerate}
\end{proposition}

\begin{proof}
By the variational definition of \(\lambda_1(X)\) and the homogeneity of
the Rayleigh quotient, for each \(j\geq1\) we may choose
\(u_j\in C_c^\infty(X)\) such that \(\int_Xu_j^2\,dV=1\) and
\[
  \frac{(m-1)^2}{4}
  \leq\int_X|\nabla u_j|^2\,dV
  \leq\frac{(m-1)^2}{4}+\frac1j.
\]
Set
\[
  \varepsilon_j
  :=\int_X|\nabla u_j|^2\,dV-\frac{(m-1)^2}{4}
  \quad\text{so that}\quad
  0\leq\varepsilon_j\leq\frac1j.
\]
Thus \(\varepsilon_j\to0\), and
\eqref{eq:minimizing-sequence} follows.

Fix \(\xi_0\in\partial_\infty X\).  By
Proposition~\ref{prop:busemann-comparison},
\[
  B_{\xi_0}\in C^2(X),\quad
  |\nabla B_{\xi_0}|=1,\quad
  \Delta B_{\xi_0}\geq m-1.
\]
Applying Lemma~\ref{lem:mckean-square} with \(B=B_{\xi_0}\) gives the exact
decomposition
\[
  \varepsilon_j
  =\int_X\left|\nabla u_j+\frac{m-1}{2}u_j\nabla B_{\xi_0}\right|^2\,dV
   +\frac{m-1}{2}\int_X
    \bigl(\Delta B_{\xi_0}-(m-1)\bigr)u_j^2\,dV.
\]
Proposition~\ref{prop:busemann-comparison} also shows that both terms on the
right are nonnegative.  Hence
\begin{align*}
  0\leq
  \int_X\left|
    \nabla u_j+\frac{m-1}{2}u_j\nabla B_{\xi_0}
  \right|^2\,dV
  &\leq\varepsilon_j\longrightarrow0,\\
  0\leq
  \int_X\bigl(\Delta B_{\xi_0}-(m-1)\bigr)u_j^2\,dV
  &\leq\frac{2\varepsilon_j}{m-1}\longrightarrow0.
\end{align*}
These are precisely \eqref{eq:first-defect} and
\eqref{eq:second-defect}.
\end{proof}

The two vanishing statements play different roles below.  The second defect
will place the limiting measure on the zero set of the Busemann defect,
whereas the first will yield stationarity and then leaf saturation of its
support.  Since the measures \(u_j^2dV\) may escape to infinity on \(X\), we
now pass to the compact horospherical suspension, where their total mass is
preserved.

\section{The horospherical suspension}
\label{sec:suspension}

We return to the cocompact setting of Theorem~\ref{thm:main}.  Let
\((M^m,g)\) be a closed Riemannian manifold with \(\sec_g\leq-1\), and let
\[
  \pi\colon(\wti M,\wti g)\longrightarrow(M,g)
\]
be its universal Riemannian covering.  Since \(M\) is compact, its sectional
curvature is also bounded below.  Thus \((\wti M,\wti g)\) satisfies the
curvature bounds in \eqref{eq:pinched-curvature}, and all the results of
Section~\ref{sec:busemann-defects} apply with \(X=\wti M\).

The next step is to extract a nonzero limit from the probability measures
\(u_j^2\,dV\).  On the noncompact manifold \(\wti M\), their mass may
escape to infinity.  Passing to the compact quotient \(M\) would prevent
this loss, but it would forget the boundary point \(\xi_0\) that defines
the Busemann function and its defect.  This boundary point cannot be held
fixed while the spatial variable is recentered: a deck transformation
\(\gamma\) sends the pair \((x,\xi_0)\) to
\((\gamma x,\gamma\xi_0)\).

The compact space must therefore include both variables and identify them
equivariantly.  The horospherical suspension defined below has these
properties.  The map
\[
  x\longmapsto[x,\xi_0]
\]
pushes \(u_j^2\,dV\) to a probability measure on the compact suspension,
while preserving the Busemann information needed in the equality
argument.

\subsection{The suspension and the unit tangent bundle}

Identify \(\Gamma=\pi_1(M)\) with the deck group of
\(\pi\colon\wti M\to M\).  Each \(\gamma\in\Gamma\) is an isometry of
\(\wti M\) and therefore acts on \(\partial_\infty\wti M\).  We use the
diagonal action
\[
  \gamma\cdot(x,\xi):=(\gamma x,\gamma\xi).
\]
The \emph{horospherical suspension} is the quotient
\begin{equation}\label{eq:suspension}
  Z
  :=
  (\wti M\times\partial_\infty\wti M)/\Gamma.
\end{equation}
The quotient is equipped with the quotient topology, and the class of
\((x,\xi)\) is denoted by \([x,\xi]\).  This is the visual-boundary form
of Ledrappier's horospherical suspension (see
\cite{Ledrappier2010}, Section~1, and
\cite{LedrappierWang2010}, Section~2).

We first identify the space before taking the quotient.  For
\((x,\xi)\in\wti M\times\partial_\infty\wti M\), let \(c_{x,\xi}\) be
the unique unit-speed geodesic ray from \(x\) to \(\xi\), and let
\[
  v_{x,\xi}:=\dot c_{x,\xi}(0)\in S_x\wti M.
\]
Define
\[
  \Phi\colon
  \wti M\times\partial_\infty\wti M
  \longrightarrow S\wti M,
  \quad
  \Phi(x,\xi):=v_{x,\xi}.
\]
Every unit vector determines its base point and its forward endpoint, so
\(\Phi\) is bijective.  Its inverse is
\[
  \Phi^{-1}(v)
  =
  \bigl(p(v),c_v(+\infty)\bigr),
\]
where \(p\colon S\wti M\to\wti M\) is the bundle projection and
\(c_v(t):=\exp_{p(v)}(tv)\).  Proposition~\ref{prop:busemann-comparison}
\textup{(iii)} gives the continuity of
\((x,\xi)\mapsto v_{x,\xi}\), while the cone topology gives the
continuity of \(v\mapsto c_v(+\infty)\).  Hence \(\Phi\) is a
homeomorphism.

The deck group acts on \(S\wti M\) by
\[
  \gamma\cdot v:=d\gamma_x(v),
  \quad v\in S_x\wti M.
\]
Since \(\gamma\) is an isometry, \(\gamma\circ c_{x,\xi}\) is the
unit-speed ray from \(\gamma x\) to \(\gamma\xi\).  Uniqueness gives
\[
  \gamma\circ c_{x,\xi}=c_{\gamma x,\gamma\xi}.
\]
Differentiating at \(t=0\) gives
\[
  v_{\gamma x,\gamma\xi}
  =d\gamma_x(v_{x,\xi}).
\]
Thus \(\Phi\) is \(\Gamma\)-equivariant and induces a homeomorphism
\[
  \Phi_\Gamma\colon
  Z\longrightarrow S\wti M/\Gamma,
  \quad
  \Phi_\Gamma([x,\xi])=[v_{x,\xi}].
\]

We next identify the quotient \(S\wti M/\Gamma\).  Because \(\pi\) is a
local isometry, its differential maps unit vectors to unit vectors and
defines
\[
  \Psi\colon S\wti M/\Gamma\longrightarrow SM,
  \quad
  \Psi([v])=d\pi_x(v),
  \quad v\in S_x\wti M.
\]
This map is well defined: since \(\pi\circ\gamma=\pi\),
\[
  d\pi_{\gamma x}\circ d\gamma_x=d\pi_x.
\]
Every unit vector in \(SM\) lifts to a unit vector in \(S\wti M\), and
any two lifts differ by a deck transformation.  Thus \(\Psi\) is
bijective.  In evenly covered neighborhoods, both \(\Psi\) and its
inverse are represented by the differential of a local isometry, so
\(\Psi\) is a homeomorphism.

Composing the two quotient maps gives
\begin{equation}\label{eq:suspension-unit-tangent}
  \overline\Phi
  :=
  \Psi\circ\Phi_\Gamma
  \colon Z\longrightarrow SM,
  \quad
  \overline\Phi([x,\xi])
  :=
  d\pi_x(v_{x,\xi}).
\end{equation}
Equivalently, the formula is independent of the representative because
\[
  d\pi_{\gamma x}(v_{\gamma x,\gamma\xi})
  =
  d\pi_{\gamma x}\bigl(d\gamma_x(v_{x,\xi})\bigr)
  =
  d\pi_x(v_{x,\xi}).
\]
Both factors in the definition of \(\overline\Phi\) are homeomorphisms;
hence
\[
  Z\cong SM.
\]
Since \(M\) is closed, \(SM\), and hence \(Z\), is compact and metrizable.

\subsection{The suspension as a foliated space}
\label{subsec:suspension-foliated-space}

We begin with the regularity convention used below.  Let
\(r\in\mathbb N\cup\{\infty\}\).  An \(m\)-dimensional
\(C^{r,0}\) \emph{foliated space} is a topological space \(X\) equipped
with charts
\[
  \psi_\alpha\colon U_\alpha\longrightarrow B_\alpha\times T_\alpha,
\]
where \(B_\alpha\subset\mathbb R^m\) is an open ball and \(T_\alpha\) is
a topological space, such that every coordinate change has the form
\[
  (x,\tau)\longmapsto
  \bigl(h_{\beta\alpha}(x,\tau),
        \theta_{\beta\alpha}(\tau)\bigr).
\]
For each fixed \(\tau\), the map
\(x\mapsto h_{\beta\alpha}(x,\tau)\) is \(C^r\), and all its
\(x\)-derivatives through order \(r\) depend continuously on
\((x,\tau)\).  When \(r=\infty\), this condition is imposed at every
finite order.  Thus \(x\) is the differentiable \emph{leaf variable},
whereas \(\tau\) is only a continuous \emph{transverse variable}.

\begin{definition}
\label{def:Ck0}
Let \(k\) be a nonnegative integer, with \(k\leq r\) when
\(r<\infty\).  A function \(f\colon X\to\mathbb R\) is of class
\(C^{k,0}\) if, in every foliated chart, all derivatives
\[
  \partial_x^I(f\circ\psi_\alpha^{-1})(x,\tau),
  \quad |I|\leq k,
\]
exist and are continuous in both variables.  No transverse derivative is
required.  For a leafwise tensor, the same condition is imposed on its
coordinate components.
\end{definition}

The sets
\[
  \psi_\alpha^{-1}(B_\alpha\times\{\tau\}),
  \quad \tau\in T_\alpha,
\]
are the \emph{plaques}.  A \emph{leaf} is a maximal subset connected by
finite chains of intersecting plaques.  The leaf coordinates make every
leaf a \(C^r\) \(m\)-manifold, although its manifold topology may be
finer than the subspace topology inherited from \(X\).  Unless stated
otherwise, every differential operator below acts only along the leaves.

Let
\[
  q\colon
  \wti M\times\partial_\infty\wti M
  \longrightarrow Z
\]
be the quotient map.  Since the diagonal action sends
\[
  \gamma\bigl(\wti M\times\{\xi\}\bigr)
  =
  \wti M\times\{\gamma\xi\},
\]
the image of a slice depends only on the \(\Gamma\)-orbit of \(\xi\).
Set
\[
  \mc L_\xi
  :=q\bigl(\wti M\times\{\xi\}\bigr)
  =\{[x,\xi]:x\in\wti M\},
\]
and define
\[
  q_\xi\colon\wti M\longrightarrow\mc L_\xi,
  \quad
  q_\xi(x):=[x,\xi].
\]
In Ledrappier's notation, \(\mc L_\xi\) is the leaf \(W_\xi\) of the
Busemann lamination; see \cite{Ledrappier2010}, Section~1.
If \(\mc L_\xi\cap\mc L_\eta\neq\varnothing\), then
\([x,\xi]=[y,\eta]\) for some \(x,y\in\wti M\), so
\((y,\eta)=(\gamma x,\gamma\xi)\) for some \(\gamma\in\Gamma\).  Hence
\(\eta=\gamma\xi\).  The converse follows directly from the diagonal
action.  Therefore
\[
  \mc L_\xi=\mc L_\eta
  \quad\Longleftrightarrow\quad
  \eta\in\Gamma\xi,
\]
and consequently
\[
  \mc F:=\{\mc L_\xi:\xi\in\partial_\infty\wti M\}
\]
is a partition of \(Z\).

\begin{lemma}
\label{lem:suspension-foliation}
The quotient \(Z\) is a compact \(m\)-dimensional
\(C^{\infty,0}\) foliated space, and its leaves are the sets
\(\mc L_\xi\).  For every \(\xi\in\partial_\infty\wti M\), the map
\(q_\xi\colon\wti M\to\mc L_\xi\) is a Riemannian covering with deck
group
\(\Gamma_\xi=\{\gamma\in\Gamma:\gamma\xi=\xi\}\), and hence induces a
Riemannian isometry
\(\wti M/\Gamma_\xi\to\mc L_\xi\).  In particular, every leaf
\(\mc L_\xi\) is a smooth, connected, complete Riemannian manifold.
\end{lemma}

\begin{proof}
\emph{Foliated atlas.}
Let \(U\subset M\) be an evenly covered coordinate ball, and choose one
lift \(\wti U\subset\wti M\).  Set
\[
  Z_{\wti U}:=q\bigl(\wti U\times\partial_\infty\wti M\bigr).
\]
The quotient map \(q\) is open, and distinct deck translates of
\(\wti U\) are disjoint.  Hence \(Z_{\wti U}\) is open in \(Z\), and
\(q\) restricts to a homeomorphism from
\(\wti U\times\partial_\infty\wti M\) onto \(Z_{\wti U}\).  Choose
smooth coordinates
\(\varphi_U\colon\wti U\to B_U\subset\mathbb R^m\), and set
\[
  \psi_{\wti U}\colon Z_{\wti U}
  \longrightarrow B_U\times\partial_\infty\wti M,
  \quad
  \psi_{\wti U}([x,\xi])
  :=
  \bigl(\varphi_U(x),\xi\bigr).
\]
This is a homeomorphism.  As \(U\) ranges over an evenly covered
coordinate cover of \(M\), these charts cover \(Z\), since every class
has a representative in one of the chosen lifts.  Their plaques are the
sets \(q(\wti U\times\{\xi\})\), with
\(\xi\in\partial_\infty\wti M\).

An overlap of two such charts decomposes into open pieces on each of
which the two representatives are related by a fixed
\(\gamma\in\Gamma\).  On such a piece the coordinate change is
\[
  (x,\xi)
  \longmapsto
  \left(
    \varphi_V\circ\gamma\circ\varphi_U^{-1}(x),
    \gamma\xi
  \right).
\]
The first component is smooth and independent of \(\xi\), while the
second is the boundary homeomorphism
\(\xi\mapsto\gamma\xi\).  Thus these charts form a
\(C^{\infty,0}\) foliated atlas of leaf dimension \(m\).

\smallskip
\noindent\emph{Identification of the leaves.}
In the chart determined by \(\wti U\), write
\(P(\wti U,\eta):=q(\wti U\times\{\eta\})\subset\mc L_\eta\).
If \(P(\wti U,\eta)\) and \(P(\wti V,\zeta)\) intersect, then
\([u,\eta]=[v,\zeta]\) for some \(u\in\wti U\) and \(v\in\wti V\).
It follows that
\((v,\zeta)=(\gamma u,\gamma\eta)\) for some \(\gamma\in\Gamma\).
In particular, \(\zeta=\gamma\eta\), and hence
\(\mc L_\eta=\mc L_\zeta\).  Therefore every plaque chain starting at
\([x,\xi]\) is contained in \(\mc L_\xi\).

Conversely, let \([y,\xi]\in\mc L_\xi\), and choose a path
\(c\colon[0,1]\to\wti M\) from \(x\) to \(y\).  For every
\(t_0\in[0,1]\), there are an interval \(I_{t_0}\), a chart
\(Z_{\wti U}\), and \(\gamma\in\Gamma\) such that
\(\gamma c(I_{t_0})\subset\wti U\).  In that chart the transverse
coordinate of \(q_\xi\circ c\) on \(I_{t_0}\) is the fixed point
\(\gamma\xi\), so this part of the path lies in one plaque.  Compactness
of \([0,1]\) gives a finite subdivision into such subarcs.  Their plaques
meet successively and form a chain from \([x,\xi]\) to \([y,\xi]\).
Therefore
\(\operatorname{Leaf}([x,\xi])=\mc L_\xi\).

\smallskip
\noindent\emph{Leaf metrics and completeness.}
Use the chosen lift \(\wti U\) to give each plaque the restriction of
\(\wti g\).  This definition is independent of the chart because every
leafwise coordinate change is induced by a deck isometry.  It therefore
defines a smooth Riemannian metric on every leaf, and
\(q_\xi\colon\wti M\to\mc L_\xi\) is a local isometry.  Moreover,
\[
  q_\xi(x)=q_\xi(y)
  \quad\Longleftrightarrow\quad
  y=\gamma x
  \quad\text{for some }\gamma\in
  \Gamma_\xi:=\{\gamma\in\Gamma:\gamma\xi=\xi\}.
\]
The subgroup \(\Gamma_\xi\) acts freely and properly discontinuously by
isometries.  Hence \(\wti M\to\wti M/\Gamma_\xi\) is a Riemannian
covering.  The displayed description of the fibers shows that \(q_\xi\)
factors through a bijective local isometry
\(\overline q_\xi\colon\wti M/\Gamma_\xi\to\mc L_\xi\), which is
therefore a Riemannian isometry.  Consequently, \(q_\xi\) is a
Riemannian covering with deck group \(\Gamma_\xi\).
The leaf \(\mc L_\xi\) is connected because \(\wti M\) is connected,
and it is complete because each of its geodesics lifts to the complete
manifold \((\wti M,\wti g)\).

Finally, \eqref{eq:suspension-unit-tangent} identifies \(Z\) with \(SM\),
which is compact because \(M\) is closed.
\end{proof}

\begin{unnumberedremark}
A leaf should not be confused with the base manifold \(M\).  The
universal covering induces a well-defined projection
\[
  \Pi\colon Z\longrightarrow M,
  \quad
  \Pi([x,\eta])=\pi(x).
\]
Under \(Z\cong SM\), this is the footpoint projection.  Its restriction
\(\Pi_\xi\colon\mc L_\xi\to M\) is the covering associated with the
subgroup \(\Gamma_\xi\subseteq\Gamma\), and
\[
  \Pi_\xi\circ q_\xi=\pi,
  \qquad
  \mc L_\xi\cong\wti M/\Gamma_\xi,
  \qquad
  M=\wti M/\Gamma.
\]

Here \(\Gamma\) is a torsion-free non-elementary word-hyperbolic group.
The stabilizer of a boundary point satisfies
\[
  \Gamma_\xi=\{e\}
  \quad\text{or}\quad
  \Gamma_\xi\cong\mathbb Z
\]
(see \cite{Bowditch1998}).  Since \(\Gamma\) is not virtually
cyclic, \(\Gamma_\xi\) has infinite index.  Thus \(\Pi_\xi\) is
infinite-sheeted and every leaf is noncompact: it is isometric to
\(\wti M\) when \(\Gamma_\xi=\{e\}\), and otherwise to a cyclic quotient
of \(\wti M\).

This does not conflict with the compactness of \(Z\).  The intrinsic
topology of a leaf may be finer than its subspace topology, so a
noncompact leaf may be nonproper or even dense in \(Z\), just as an
irrational line is dense in a torus.  Intrinsically, however, every leaf
is smooth, connected, and complete, and it inherits the local curvature
bounds of \(\wti M\).  These are precisely the properties used in the
leafwise elliptic and heat-kernel arguments below.
\end{unnumberedremark}


\begin{proposition}
\label{prop:busemann-field-regularity}
Let \((M^m,g)\) be a closed Riemannian manifold with \(\sec_g\leq-1\), and
let
\[
  Z=(\wti M\times\partial_\infty\wti M)/\Gamma.
\]
The Busemann vector field on the covering space,
\[
  \widetilde V(x,\xi):=\nabla B_\xi(x)
\]
is \(\Gamma\)-equivariant and therefore descends to a leafwise unit vector
field \(V\) on \(Z\).  The maps
\[
  (x,\xi)\longmapsto\widetilde V(x,\xi),
  \quad
  (x,\xi)\longmapsto\nabla_x\widetilde V(x,\xi)
\]
are continuous on \(\wti M\times\partial_\infty\wti M\).  Moreover, for
each fixed \(\xi\), the map \(x\mapsto\widetilde V(x,\xi)\) is smooth.
Consequently,
\[
  V\in C^{1,0}(T\mc F),
  \qquad
  V|_{\mc L_\xi}\in C^\infty(T\mc L_\xi)
  \quad\text{for every }\xi.
\]
\end{proposition}

\begin{proof}
\emph{Equivariance and unit length.}
For \(\gamma\in\Gamma\), the Busemann functions
\(B_{\gamma\xi}\circ\gamma\) and \(B_\xi\) differ by a constant.  Since
\(\gamma\) is an isometry, taking gradients gives
\[
  \nabla B_{\gamma\xi}(\gamma x)
  =d\gamma_x\bigl(\nabla B_\xi(x)\bigr).
\]
Thus \(\widetilde V\) is equivariant and defines \(V\) on the quotient.
Equation~\eqref{eq:unit-busemann-gradient} gives \(|V|=1\).

\smallskip
\noindent\emph{Joint continuity to first leafwise order.}
Since \(M\) is compact, there is \(A\geq1\) such that
\[
  -A^2\leq\sec_g\leq-1.
\]
Proposition~\ref{prop:busemann-comparison}\textup{(iii)} therefore
applies to \(\wti M\) and gives joint continuity of \(\widetilde V\)
and its covariant derivative in the \(x\)-variable:
\[
  \nabla_x\widetilde V=(\nabla^2B_\xi)^\sharp.
\]
In a flow box from Lemma~\ref{lem:suspension-foliation}, let \(V^i\)
denote the leafwise coordinate components.  Then
\[
  \partial_jV^i
  =
  (\nabla_jV)^i-\Gamma^i_{jk}V^k.
\]
The Christoffel symbols in these lifted coordinates are smooth in the
leaf variable and independent of the boundary coordinate.  Hence both
\(V^i\) and \(\partial_jV^i\) are jointly continuous, which is precisely
\(V\in C^{1,0}(T\mc F)\).

\smallskip
\noindent\emph{Smoothness for fixed \(\xi\).}
Before taking the quotient, \(\Phi\) sends the slice
\(\wti M\times\{\xi\}\) onto
\[
  W^{ws}(\xi)
  =
  \{v_{x,\xi}:x\in\wti M\}
  \subset S\wti M.
\]
Two unit vectors lie in the same weak-stable leaf exactly when their
forward geodesic rays have the same endpoint at infinity.  The adjective
``weak'' indicates that a shift in the geodesic-flow parameter is
allowed; fixing also the Busemann level gives the strong-stable leaves.
Under \(\overline\Phi\colon Z\to SM\), the leaf \(\mc L_\xi\)
corresponds to the immersed weak-stable leaf
\(W^{ws}(\xi)/\Gamma_\xi\).

The geodesic flow on \(SM\) is a smooth Anosov flow
(see \cite{Anosov1967}).  By the stable manifold theorem, its weak-stable
leaves are smooth immersed submanifolds (see
\cite{HirschPughShub1977}, Theorem~4.1\textup{(e)} and the remarks
following it); their transverse dependence is generally only H\"older.
The lifted weak-stable foliation on \(S\wti M\) has the same leafwise
regularity.  Since there is a unique ray from each \(x\) to \(\xi\), the
footpoint projection
\(p\colon S\wti M\to\wti M\) restricts to a bijection
\[
  p|_{W^{ws}(\xi)}\colon W^{ws}(\xi)\longrightarrow\wti M.
\]

To prove that this bijection is a diffeomorphism, fix
\(v=v_{x,\xi}\).  The tangent space of the weak-stable leaf is
\[
  T_vW^{ws}(\xi)=\mathbb RX(v)\oplus E^s(v),
\]
where \(X\) is the geodesic vector field, and \(dp_v(X(v))=v\).
A vector \(z\in E^s(v)\) is represented by a normal Jacobi field \(J_z\)
with \(dp_v(z)=J_z(0)\) and \(J_z(t)\to0\).  If \(J_z(0)=0\), then the
Jacobi equation and \(\sec_g\leq-1\) show that
\(\frac12|J_z|^2\) is strictly convex unless \(J_z\equiv0\).  Since its
value and first derivative vanish at \(t=0\), strict convexity would make
it increase and prevent \(J_z(t)\to0\).  Thus \(J_z\equiv0\).
Furthermore, \(dp_v(E^s(v))\perp v\), whereas \(dp_v(X(v))=v\).
Consequently, \(dp_v\) is injective on
\(\mathbb RX(v)\oplus E^s(v)\), and hence is an isomorphism by dimension.
The inverse function theorem now shows that
\(p|_{W^{ws}(\xi)}\) is a local diffeomorphism.  Since it is also
bijective, it is a diffeomorphism.  Its inverse
\[
  \wti M\longrightarrow W^{ws}(\xi),
  \quad
  x\longmapsto v_{x,\xi},
\]
is therefore smooth.  Finally, \eqref{eq:unit-busemann-gradient} gives
\[
  \nabla B_\xi(x)=-v_{x,\xi},
\]
so \(x\mapsto\widetilde V(x,\xi)\) is smooth.  The leaf charts from
Lemma~\ref{lem:suspension-foliation} then show that
\(V|_{\mc L_\xi}\) is smooth.
\end{proof}

\begin{unnumberedremark}
The two regularity statements serve different purposes.
\(V\in C^{1,0}\) means that \(V\) and its first leafwise derivative vary
continuously in the transverse parameter; it does not require any
transverse derivative.  Smoothness of \(V|_{\mc L_\xi}\) gives
derivatives of every order on a fixed leaf, without asserting that the
higher derivatives vary continuously with \(\xi\).  The first property
provides continuous coefficients for the global leafwise diffusion; the
second permits classical heat-kernel arguments on an individual leaf.
\end{unnumberedremark}

\subsection{Leafwise differential calculus}

All derivatives in this subsection are taken only in the leaf variable.
For \(\xi\in\partial_\infty\wti M\), let
\[
  \iota_\xi\colon\wti M\longrightarrow\mc L_\xi,
  \quad \iota_\xi(x)=[x,\xi].
\]
Lemma~\ref{lem:suspension-foliation} identifies this map with the
Riemannian covering
\[
  \wti M\longrightarrow
  \wti M/\Gamma_\xi\cong\mc L_\xi.
\]

\smallskip
\noindent\emph{Leafwise regularity.}
For a function \(F\colon Z\to\mathbb R\), set
\[
  F_\xi:=F\circ\iota_\xi\colon\wti M\longrightarrow\mathbb R.
\]
For an integer \(k\geq0\), we say that \(F\) is \(C^k\)
\emph{along the leaves} if \(F|_{\mc L}\) is \(C^k\) for every leaf
\(\mc L\).  Since \(\iota_\xi\) is a covering, this is equivalent to
\[
  F_\xi\in C^k(\wti M)
  \quad
  \text{for every }\xi\in\partial_\infty\wti M.
\]
This condition concerns one leaf at a time and imposes no continuity as
the leaf varies.

The notation \(F\in C^{k,0}(Z)\) is understood in the sense of
Definition~\ref{def:Ck0}.  It is stronger than being \(C^k\) along
each leaf separately, because the leafwise derivatives must also vary
continuously in the transverse parameter.  As before, no transverse
derivative is required.

\smallskip
\noindent\emph{Leafwise differential operators.}
If \(F\) is \(C^1\) along the leaves, the leafwise metric defines
\(\nabla^{\mathrm{leaf}}F\) by restricting \(F\) to each leaf.  If
\(F\) is \(C^2\) along the leaves, define
\(\Delta^{\mathrm{leaf}}F\) in the same way.  The covering
\(\iota_\xi\) is a local isometry, so the naturality of the gradient and
Laplacian gives
\begin{equation}\label{eq:leafwise-lift-calculus}
\begin{aligned}
  d(\iota_\xi)_x\bigl(\nabla F_\xi(x)\bigr)
  &=\nabla^{\mathrm{leaf}}F(\iota_\xi(x)),\\
  \Delta F_\xi
  &=(\Delta^{\mathrm{leaf}}F)\circ\iota_\xi,
\end{aligned}
\end{equation}
where the unadorned \(\nabla\) and \(\Delta\) are the Riemannian operators on
\((\wti M,\wti g)\).  If \(dV_{\mc L_\xi}\) denotes the Riemannian
volume density of the leaf, then
\[
  \iota_\xi^*dV_{\mc L_\xi}=dV_{\wti g}.
\]
Consequently, a pointwise leafwise differential identity may be checked
after pulling it back to \(\wti M\) with \(\xi\) fixed.  Integration by
parts may likewise be carried out in plaque coordinates, or intrinsically
on a leaf under the usual compact-support or integrability hypotheses.

\section{The Busemann defect and the limiting measure}
\label{sec:limiting-measure}

Pushing the quadratic mass of a minimizing sequence to the compact
suspension \(Z\) prevents loss of mass.  The Busemann defect descends to
\(Z\), and the limiting measure is supported in its zero set.

Define the \emph{Busemann defect} on the covering space by
\[
  \widetilde D(x,\xi)=\Delta B_\xi(x)-(m-1),
  \quad (x,\xi)\in\wti M\times\partial_\infty\wti M.
\]
By Proposition~\ref{prop:busemann-comparison}, \(\widetilde D\) is
continuous and nonnegative.

For \(\gamma\in\Gamma\), the functions
\(B_{\gamma\xi}\circ\gamma\) and \(B_\xi\) differ by a constant.  Since
\(\gamma\) is an isometry,
\[
  \Delta B_{\gamma\xi}(\gamma x)=\Delta B_\xi(x).
\]
Thus \(\widetilde D\) is \(\Gamma\)-invariant and descends to the
continuous function
\begin{equation}\label{eq:busemann-defect}
  D\colon Z\longrightarrow[0,\infty),
  \quad
  D([x,\xi])=\Delta B_\xi(x)-(m-1)
\end{equation}
for every \((x,\xi)\in\wti M\times\partial_\infty\wti M\).  The set
\(D^{-1}(0)\) consists of the classes represented by pairs
\((x,\xi)\) for which \(\Delta B_\xi(x)=m-1\).

To construct the limiting measure, assume
\(\lambda_1(\wti M)=(m-1)^2/4\), and fix
\(\xi_0\in\partial_\infty\wti M\).  Choose a normalized minimizing sequence
\(u_j\) as in Proposition~\ref{prop:defects}, with \(X=\wti M\).  The
measure \(u_j^2dV\) is a probability measure on \(\wti M\).  Push it to
the compact suspension under the continuous map
\[
  \iota_{\xi_0}\colon\wti M\longrightarrow Z,
  \quad x\longmapsto[x,\xi_0],
\]
and define
\begin{equation}\label{eq:mu-j}
  \mu_j=(\iota_{\xi_0})_*(u_j^2dV).
\end{equation}
Each \(\mu_j\) is a probability measure, and for \(f\in C(Z)\),
\begin{equation}\label{eq:mu-j-testing}
  \int_Z f\,d\mu_j
  =\int_{\wti M}f([x,\xi_0])\,u_j(x)^2\,dV(x).
\end{equation}

\begin{proposition}
\label{prop:defect-support}
After passing to a subsequence, \(\mu_{j_k}\rightharpoonup\mu\), where
\(\mu\) is a probability measure satisfying
\[
  \int_ZD\,d\mu=0,
  \qquad
  \supp\mu\subseteq D^{-1}(0).
\]
\end{proposition}

\begin{proof}
The probability measures on the compact metrizable space \(Z\) are
weakly sequentially compact (see \cite{Bogachev2007},
Theorem~8.6.7).  Thus some subsequence converges weakly to a probability
measure \(\mu\).  Continuity of \(D\), weak convergence, the pushforward
formula \eqref{eq:mu-j-testing}, and the definition
\eqref{eq:busemann-defect} give
\[
  \int_ZD\,d\mu
  =\lim_{k\to\infty}\int_ZD\,d\mu_{j_k}
  =\lim_{k\to\infty}
    \int_{\wti M}\bigl(\Delta B_{\xi_0}-(m-1)\bigr)u_{j_k}^2\,dV
  =0.
\]
The last equality is \eqref{eq:second-defect}.  Since \(D\geq0\),
the identity forces \(D=0\) on \(\supp\mu\), which proves the stated
support containment.
\end{proof}

\section{Leafwise stationarity and support saturation}
\label{sec:stationarity}

Fix a weak subsequential limit \(\mu\) as in
Proposition~\ref{prop:defect-support}.  The first defect gives
infinitesimal stationarity for a leafwise drift operator.  The general
results of Section~\ref{sec:candel-semigroup-construction} then imply
stationarity and leaf saturation of \(\supp\mu\).

\subsection{The leafwise drift operator}

Let \(V\) be the Busemann field from
Proposition~\ref{prop:busemann-field-regularity}.  It is \(C^{1,0}\),
smooth on each leaf, and its lift to the covering
\(\iota_\xi\colon\wti M\to\mc L_\xi\) is \(\nabla B_\xi\).

Let \(f\in C(Z)\), and assume that \(f|_{\mc L}\) is \(C^2\) on every leaf
\(\mc L\).  Define
\begin{equation}\label{eq:leafwise-L}
  Lf
  :=\Delta^{\mathrm{leaf}}f
    -(m-1)\langle V,\nabla^{\mathrm{leaf}}f\rangle.
\end{equation}
By \eqref{eq:leafwise-lift-calculus}, its pullback to \(\wti M\) is
\begin{equation}\label{eq:leafwise-L-lift}
  (Lf)\circ\iota_\xi
  =\Delta(f\circ\iota_\xi)
   -(m-1)\bigl\langle\nabla B_\xi,
                 \nabla(f\circ\iota_\xi)\bigr\rangle.
\end{equation}
Set \(A:=L/2\).  We use its classical continuous domain
\begin{equation}\label{eq:generator-test-domain}
  \mc D_0(A)
  :=\left\{f\in C(Z):
    f|_{\mc L}\in C^2(\mc L)\text{ for every leaf }\mc L,
    \quad Af\in C(Z)
  \right\}.
\end{equation}
Since \(A=L/2\), the condition \(Af\in C(Z)\) is equivalent to
\(Lf\in C(Z)\).  This is the classical domain used in
Section~\ref{sec:candel-semigroup-construction} and also the test domain
in Suzaki's definition of an \(A\)-harmonic measure (see
\cite{Suzaki2015}, p.~256).

\begin{lemma}
\label{lem:leafwise-L-elliptic}
In every quotient product chart, the principal coefficients of \(L\) are
of class \(C^{\infty,0}\), its first-order coefficients are of
class \(C^{1,0}\), and all coefficients are smooth on each plaque.  Its
principal symbol is \((g^{\mathrm{leaf}})^{-1}\), and it has no
zeroth-order term.  In particular, \(L\) is uniformly elliptic along the
leaves and \(L1=0\).
\end{lemma}

\begin{proof}
In a foliated chart \(B\times T\), write
\[
  g^{\mathrm{leaf}}=g_{ij}(y,\tau)\,dy^i\,dy^j,
  \quad
  V=V^j(y,\tau)\partial_j.
\]
The standard coordinate formula for the Laplace--Beltrami operator gives
\[
  Lf
  =g^{ij}\partial_i\partial_jf
   +\left[
      \frac1{\sqrt{\det g}}\,
      \partial_i\bigl(\sqrt{\det g}\,g^{ij}\bigr)
      -(m-1)V^j
    \right]\partial_jf.
\]
The leafwise metric is of class \(C^{\infty,0}\), while
Proposition~\ref{prop:busemann-field-regularity} gives
\(V\in C^{1,0}\) and shows that \(V\) is smooth on each plaque.  Hence
\(g^{ij}\in C^{\infty,0}\), the displayed first-order coefficients belong
to \(C^{1,0}\), and all coefficients are smooth on each plaque.  The
principal matrix is \((g^{ij})\), so the principal symbol is
\((g^{\mathrm{leaf}})^{-1}\).  Its positive definiteness and the
compactness of \(Z\) give, after choosing a finite foliated atlas, a
constant \(\theta\in(0,1]\) such that
\[
  \theta|\zeta|^2
  \leq g^{ij}(y,\tau)\zeta_i\zeta_j
  \leq\theta^{-1}|\zeta|^2.
\]
Finally, the displayed formula contains no zeroth-order term, and hence
\(L1=0\).
\end{proof}

Lemma~\ref{lem:classical-gradient-control}, applied with
\(\ms A=A\), shows that every \(f\in\mc D_0(A)\) has a continuous,
uniformly bounded leafwise gradient.  This is the only domain
regularity needed in the first-defect argument.

\subsection{The leafwise diffusion}

The general results of Section~\ref{sec:candel-semigroup-construction}
apply to \(A=L/2\).  We state their consequences for the suspension.

\begin{lemma}
\label{lem:saturated-support}
The maximal continuous realization of \(A\) generates a conservative
Feller Markov semigroup \(\{P_t\}_{t\geq0}\) on \(C(Z)\).  Moreover:
\begin{enumerate}[label=\textup{(\roman*)}]
\item On every leaf, \(P_t\) is the intrinsic minimal heat semigroup of
\(A\).  Its heat kernel is smooth, strictly positive, and conservative
with respect to the principal metric \(g_A=2g^{\mathrm{leaf}}\).

\item A Borel probability measure \(\nu\) on \(Z\) satisfies
\[
  \int_ZP_tf\,d\nu=\int_Zf\,d\nu
  \quad(f\in C(Z),\ t\geq0)
\]
if and only if
\begin{equation}\label{eq:infinitesimal-stationarity}
  \int_ZLf\,d\nu=0
  \quad\text{for every }f\in\mc D_0(A).
\end{equation}

\item If \(\nu\) satisfies the equivalent conditions in
\textup{(ii)}, then
\[
  z\in\supp\nu\quad\Longrightarrow\quad
  \mc L(z)\subset\supp\nu.
\]
\end{enumerate}
\end{lemma}

\begin{proof}
By Lemma~\ref{lem:suspension-foliation}, \(Z\) is a compact
\(C^{\infty,0}\) foliated space.  Lemma~\ref{lem:leafwise-L-elliptic}
verifies all coefficient hypotheses for \(A\), and its principal matrix
is \((g^{ij}/2)\); hence \(g_A=2g^{\mathrm{leaf}}\).
Part~\textup{(i)} follows from
Theorem~\ref{thm:candel-diffusion}\textup{(ii)}--\textup{(iii)}.
Since \(L=2A\), part~\textup{(ii)} is
Proposition~\ref{prop:stationarity-equivalence}, and part~\textup{(iii)}
is Proposition~\ref{prop:stationary-support-saturated}.
\end{proof}

\subsection{Stationarity from the first defect}

We pass the first-defect identity to the weak limit.  Differentiating
\(u_j^2\) gives exactly the coefficient \(m-1\) appearing in \(L\).

\begin{proposition}
\label{prop:stationarity-defect}
Every weak subsequential limit \(\mu\) from
Proposition~\ref{prop:defect-support} satisfies
\begin{equation}\label{eq:mu-infinitesimal-stationarity}
  \int_ZLf\,d\mu=0
  \quad\text{for every }f\in\mc D_0(A).
\end{equation}
Moreover, its support is leaf-saturated; explicitly,
\[
  z\in\supp\mu\quad\Longrightarrow\quad
  \mc L(z)\subset\supp\mu.
\]
\end{proposition}

\begin{proof}
Use the notation of Section~\ref{sec:limiting-measure}, and choose a
subsequence such that \(\mu_{j_k}\rightharpoonup\mu\).  Fix
\(f\in\mc D_0(A)\), and set
\[
  f_0:=f\circ\iota_{\xi_0},
  \quad
  R_j=\nabla u_j+\frac{m-1}{2}u_j\nabla B_{\xi_0}.
\]
The local isometry
\(\iota_{\xi_0}\colon\wti M\to\mc L_{\xi_0}\) gives
\(f_0\in C^2(\wti M)\) and
\[
  |\nabla f_0(x)|
  =|\nabla^{\mathrm{leaf}}f(\iota_{\xi_0}(x))|
  \quad\text{for every }x\in\wti M.
\]
Lemma~\ref{lem:classical-gradient-control}, together with
\(g_A=2g^{\mathrm{leaf}}\), therefore gives
\[
  \|\nabla f_0\|_{L^\infty(\wti M)}
  \leq\|\nabla^{\mathrm{leaf}}f\|_{L^\infty(Z)}<\infty.
\]

By the pushforward identity \eqref{eq:mu-j-testing} and the lifted formula
\eqref{eq:leafwise-L-lift},
\[
  \int_ZLf\,d\mu_j
  =\int_{\wti M}u_j^2
   \bigl(\Delta f_0
   -(m-1)\langle\nabla B_{\xi_0},\nabla f_0\rangle\bigr)\,dV.
\]
Although \(f_0\) need not have compact support, \(u_j\) does, so integration
by parts gives
\[
  \int_{\wti M}u_j^2\Delta f_0\,dV
  =-2\int_{\wti M}u_j
       \langle\nabla u_j,\nabla f_0\rangle\,dV.
\]
Substitution yields
\begin{equation}\label{eq:first-defect-pairing}
  \int_ZLf\,d\mu_j
  =-2\int_{\wti M}u_j\langle R_j,\nabla f_0\rangle\,dV.
\end{equation}
By Cauchy--Schwarz,
\begin{align*}
  \left|\int_ZLf\,d\mu_j\right|
  &\leq2\|u_j\|_{L^2(\wti M)}
        \|R_j\|_{L^2(T\wti M)}
        \|\nabla f_0\|_{L^\infty(\wti M)}\\
  &\leq2\|\nabla^{\mathrm{leaf}}f\|_{L^\infty(Z)}
        \|R_j\|_{L^2(T\wti M)}
  \longrightarrow0.
\end{align*}
The convergence follows from \eqref{eq:first-defect}.  Since
\(Lf\in C(Z)\), weak convergence gives
\[
  \int_ZLf\,d\mu
  =\lim_{k\to\infty}\int_ZLf\,d\mu_{j_k}=0.
\]
Thus \eqref{eq:mu-infinitesimal-stationarity} holds.  By
Lemma~\ref{lem:saturated-support}\textup{(ii)}--\textup{(iii)},
\(\mu\) is stationary and \(\supp\mu\) is leaf-saturated.
\end{proof}

\subsection{From saturated support to Busemann equality}

Propositions~\ref{prop:defect-support} and
\ref{prop:stationarity-defect} show that \(\supp\mu\) lies in
\(D^{-1}(0)\) and is leaf-saturated.  Choosing a leaf contained in this
support therefore gives a single boundary point for which the Busemann
defect vanishes at every base point.

\begin{proposition}
\label{prop:equality-busemann}
Let \((M^m,g)\) be a closed connected Riemannian manifold with \(m\geq2\)
and \(\sec_g\leq-1\), and let \((\wti M,\wti g)\) be its universal cover.
Assume
\[
  \lambda_1(\wti M)=\frac{(m-1)^2}{4}.
\]
Then there exists \(\xi_*\in\partial_\infty\wti M\) such that
\(B_{\xi_*}\in C^2(\wti M)\) and
\[
  |\nabla B_{\xi_*}(x)|=1,
  \quad
  \Delta B_{\xi_*}(x)=m-1
\]
for every \(x\in\wti M\).
\end{proposition}

\begin{proof}
Since \(M\) is compact, there is a constant \(A\geq1\) such that
\[
  -A^2\leq\sec_g\leq-1.
\]
The same bounds hold on \((\wti M,\wti g)\), so
Proposition~\ref{prop:busemann-comparison} shows that every Busemann
function \(B_\xi\) is of class \(C^2\) and satisfies
\(|\nabla B_\xi|=1\).  We seek a boundary point for which the Laplacian
equality holds throughout \(\wti M\).

Proposition~\ref{prop:defect-support} gives a probability measure
\(\mu\), and Proposition~\ref{prop:stationarity-defect} applies to this
measure.  Together they give
\[
  \supp\mu\subseteq D^{-1}(0),
  \quad
  \supp\mu\text{ is leaf-saturated}.
\]

Since \(\mu\) is a probability measure, \(\supp\mu\neq\varnothing\).
Choose \(z_*=[x_*,\xi_*]\in\supp\mu\).  By the definition of the leaf
through \(z_*\) and leaf saturation, for every \(x\in\wti M\),
\[
  [x,\xi_*]\in\mc L(z_*)
  \subseteq\supp\mu
  \subseteq D^{-1}(0).
\]
Therefore, by \eqref{eq:busemann-defect},
\[
  0=D([x,\xi_*])
   =\Delta B_{\xi_*}(x)-(m-1)
  \quad\text{for every }x\in\wti M.
\]
Hence \(\Delta B_{\xi_*}=m-1\) throughout \(\wti M\).  The regularity and
unit-gradient identity were established above.
\end{proof}

It remains to convert this Busemann equality into a global geometric
description of \(\wti M\).

\section{From Busemann equality to hyperbolic rigidity}
\label{sec:rigidity}

In this section we pass from the Busemann equality obtained in
Proposition~\ref{prop:equality-busemann} to hyperbolic rigidity and complete
the proof of Theorem~\ref{thm:main}.  The next lemma contains the
Riemannian part of the argument: the upper curvature bound gives the
Busemann Hessian identity, integrating the gradient flow gives a
warped-product splitting, and the lower curvature bound forces the
reference horosphere to be flat.  We then apply the lemma to the universal
cover of \(M\).  The final subsection explains the two distinct places
where cocompactness is used.

\subsection{Riemannian rigidity and completion of the proof}

\begin{lemma}
\label{lem:riemannian-rigidity}
Let \((X^m,g_X)\) be a complete simply connected Riemannian manifold with
\(m\geq2\), and suppose that, for some \(A\geq1\),
\[
  -A^2\leq\sec_{g_X}\leq-1.
\]
If a Busemann function \(B\in C^2(X)\) satisfies
\[
  |\nabla B|=1,
  \quad
  \Delta B=m-1,
\]
then \(B\) is smooth,
\begin{equation}\label{eq:hessian-equality}
  \nabla^2B=g_X-dB\otimes dB.
\end{equation}
Moreover, \(H=B^{-1}(0)\) is complete, simply connected, and flat, and
\begin{equation}\label{eq:warped-product}
  (X,g_X)
  \cong
  \bigl(\R\times H,dt^2+\ee^{2t}g_H\bigr).
\end{equation}
Consequently, \((X,g_X)\) is isometric to
\(\mathbb H^m(-1)\).
\end{lemma}

\begin{proof}
The lower Hessian comparison inequality
\eqref{eq:busemann-laplacian-comparison} shows that
\[
  Q:=\nabla^2B-(g_X-dB\otimes dB)
\]
is positive semidefinite.  Its trace is
\[
  \tr_{g_X}Q
  =\Delta B-\bigl(m-|\nabla B|^2\bigr)
  =(m-1)-(m-1)=0.
\]
A positive semidefinite symmetric tensor with zero trace vanishes, proving
\eqref{eq:hessian-equality}.

The tensor identity also shows that \(B\) is smooth.  Indeed, in local
coordinates,
\[
  \partial_i\partial_jB
  =\Gamma_{ij}^{k}\partial_kB+(g_X)_{ij}
   -(\partial_iB)(\partial_jB).
\]
Starting from \(B\in C^2\), this equation and induction give
\(B\in C^\infty(X)\).  If \(T=\nabla B\), then
\eqref{eq:hessian-equality} is equivalent to
\begin{equation}\label{eq:nabla-T}
  \nabla_YT=Y-\langle Y,T\rangle T
\end{equation}
for every vector field \(Y\).  In particular,
\begin{equation}\label{eq:T-geodesic}
  \nabla_TT=0,
  \quad
  \nabla_YT=Y\quad\text{when }Y\perp T.
\end{equation}
Thus the integral curves of \(T\) are unit-speed geodesics, while the level
sets of \(B\) are totally umbilical with second fundamental form equal to
their induced metric.

Set \(H=B^{-1}(0)\), with its induced metric \(g_H\).  We next integrate
the vector field \(T=\nabla B\).
Since \(|T|=1\) and \(X\) is complete, the flow \(\Phi_t\) of
\(T=\nabla B\) is defined for every \(t\in\R\).  Along each flow line,
\begin{equation}\label{eq:B-along-flow}
  \frac{d}{dt}B(\Phi_t(x))
  =\langle\nabla B,T\rangle=1,
\end{equation}
and hence \(B(\Phi_t(x))=B(x)+t\).  It follows that
\[
  \Phi\colon\R\times H\longrightarrow X,
  \quad
  \Phi(t,y)=\Phi_t(y),
\]
is a diffeomorphism with inverse
\[
  x\longmapsto\bigl(B(x),\Phi_{-B(x)}(x)\bigr).
\]
Since \(X\) is connected, \(H\) is connected.

Let \(g_t\) be the metric on \(H\) obtained by pulling back the induced
metric on \(B^{-1}(t)\) by \(\Phi_t\).  Extend tangent vector fields
\(Y,Z\) on \(H\) by the flow, so that \([T,Y]=[T,Z]=0\).  Using
\eqref{eq:T-geodesic},
\begin{align*}
  \frac{d}{dt}g_t(Y,Z)
  &=\langle\nabla_TY,Z\rangle+\langle Y,\nabla_TZ\rangle \\
  &=\langle\nabla_YT,Z\rangle+\langle Y,\nabla_ZT\rangle \\
  &=2g_t(Y,Z).
\end{align*}
Therefore \(g_t=\ee^{2t}g_H\).  Since \(T\) is unit and orthogonal to the
level sets, this proves \eqref{eq:warped-product}.

The level set \(H\) is closed in \(X\).  An intrinsic Cauchy sequence in
\(H\) is therefore an ambient Cauchy sequence; completeness of \(X\) gives
an ambient limit in \(H\), and local equivalence of the intrinsic and
ambient distances gives intrinsic convergence.  Thus \(H\) is complete.
Finally, \(X\cong\R\times H\) makes \(H\) a deformation retract of \(X\),
which proves the assertion about simple connectivity.

The sectional curvatures of \eqref{eq:warped-product} are as follows.
If \(Y\) is tangent to \(H\), then
\begin{equation}\label{eq:mixed-curvature}
  \sec_{g_X}(\partial_t,Y)=-1.
\end{equation}
If \(\sigma\subset T_yH\) is a two-plane and \(\sigma_t\) is the
corresponding plane tangent to \(\{t\}\times H\), then
\begin{equation}\label{eq:tangential-curvature}
  \sec_{g_X}(\sigma_t)
  =\ee^{-2t}\sec_{g_H}(\sigma)-1.
\end{equation}
These are the standard warped-product formulas for the warping function
\(\ee^t\).

It remains to use the lower curvature bound.  The splitting alone would
allow \(g_H\) to be nonflat and nonpositively curved.
If \(m=2\), then \(H\) is one-dimensional and hence flat.  Suppose
\(m\geq3\), and let \(\sigma\subset T_yH\) be a two-plane.
The tangential curvature formula \eqref{eq:tangential-curvature} and the
upper curvature bound imply
\[
  \ee^{-2t}\sec_{g_H}(\sigma)-1\leq-1,
\]
and hence \(\sec_{g_H}(\sigma)\leq0\).  If
\(\sec_{g_H}(\sigma)=-\kappa<0\), then
\[
  \sec_{g_X}(\sigma_t)=-\kappa\ee^{-2t}-1
  \longrightarrow-\infty
  \quad\text{as }t\longrightarrow-\infty.
\]
This contradicts the lower curvature bound.
Therefore \(\sec_{g_H}(\sigma)=0\) for every \(\sigma\), so \(g_H\) is flat.

Since \(H\) is complete, simply connected, and flat,
\((H,g_H)\cong(\R^{m-1},g_{\mr{Euc}})\).  With \(r=\ee^{-t}\),
\[
  dt^2+\ee^{2t}g_{\mr{Euc}}
  =\frac{dr^2+g_{\mr{Euc}}}{r^2},
\]
which is the upper-half-space metric on \(\mathbb H^m(-1)\).
\end{proof}

\begin{proof}[Proof of Theorem~\ref{thm:main}]
McKean's estimate gives the asserted lower bound.  Assume equality.
Since \(M\) is closed, the sectional curvature of \(g\) is bounded below;
thus, for some \(A\geq1\),
\[
  -A^2\leq\sec_{\wti g}\leq-1.
\]
Proposition~\ref{prop:equality-busemann} gives a Busemann function
\(B\in C^2(\wti M)\) satisfying
\[
  |\nabla B|=1,
  \quad
  \Delta B=m-1.
\]
Lemma~\ref{lem:riemannian-rigidity} therefore yields
\(\wti M\cong\bH^m(-1)\).

Conversely, if \(\wti M\cong\bH^m(-1)\), then
\[
  \lambda_1(\wti M)=\frac{(m-1)^2}{4}.
\]
The upper bound follows by applying radial cutoffs to
\(\exp(-(m-1)r/2)\), and the reverse inequality is McKean's estimate.
The covering projection \(\wti M\to M\) is a local isometry.  Thus
\(\wti M\cong\bH^m(-1)\) is equivalent to \(g\) having constant sectional
curvature \(-1\).
\end{proof}

\subsection{The role of cocompactness}
\label{subsec:role-cocompactness}

Closedness has two concrete consequences in the proof of
Theorem~\ref{thm:main}.  It yields a uniform lower sectional-curvature
bound, which is the hypothesis that forces the reference horosphere to be
flat in Lemma~\ref{lem:riemannian-rigidity}.  It also makes the horospherical
suspension \(Z\cong SM\) compact, so the probability measures associated
with a minimizing sequence cannot lose mass at infinity.  The two
constructions below show that these are separate issues.

\medskip
\noindent\emph{Failure without a lower curvature bound.}
For \(m\geq3\), let \((H^{m-1},g_H)\) be complete, simply connected,
nonflat, and satisfy \(\sec_{g_H}\leq0\).  The warped product
\[
  (X^m,g_X)
  =\bigl(\R\times H,dt^2+\ee^{2t}g_H\bigr)
\]
is complete and simply connected.  It has mixed sectional curvatures
\(-1\), while for a two-plane
\(\sigma\subset T_yH\),
\[
  \sec_{g_X}(\sigma_t)
  =\ee^{-2t}\sec_{g_H}(\sigma)-1.
\]
Thus \(\sec_{g_X}\leq-1\), but \(X\) is not hyperbolic.  Nevertheless,
\[
  \lambda_1(X)=\frac{(m-1)^2}{4}.
\]
McKean's estimate gives the lower bound.  For the reverse inequality, fix
\(0\neq\varphi\in C_c^\infty(H)\), put \(a=(m-1)/2\), and choose
\(\chi_R\in C_c^\infty(\mathbb R)\) supported in \((0,R+1)\), equal to one
on \([1,R]\), and with derivative bounded independently of \(R\).  Since
\(dV_X=\ee^{(m-1)t}dt\,dV_H\), the functions
\[
  u_R(t,y)=\chi_R(t)\ee^{-at}\varphi(y)
\]
satisfy
\[
  \frac{\int_X|\nabla u_R|^2\,dV_X}
       {\int_Xu_R^2\,dV_X}
  =a^2+O(R^{-1}).
\]
Indeed, the denominator is comparable to \(R\), while the cutoff terms
and the \(H\)-gradient term remain bounded.  Hence the Rayleigh quotients
converge to \(a^2\).  Negative curvature in \(g_H\) is amplified as
\(t\to-\infty\), so \(\sec_{g_X}\) has no uniform lower bound.  This is
precisely the obstruction excluded by the two-sided curvature hypothesis
in Lemma~\ref{lem:riemannian-rigidity}.

\medskip
\noindent\emph{Failure of tightness at finite volume.}
The lower curvature bound may still hold when the suspension is no longer
compact.  The next proposition is a spectral criterion on the
universal cover; the example following it constructs a finite-volume
quotient satisfying the criterion.

\begin{proposition}
\label{prop:hyperbolic-horoball-criterion}
Let \((X^m,g_X)\), \(m\geq2\), be complete and simply connected, and
suppose that
\[
  \sec_{g_X}\leq-1.
\]
If \(X\) contains an open subset isometric to a horoball in
\(\mathbb H^m(-1)\), then
\begin{equation}\label{eq:cusp-spectrum}
  \lambda_1(X)=\frac{(m-1)^2}{4}.
\end{equation}
\end{proposition}

\begin{proof}
McKean's inequality gives
\[
  \lambda_1(X)\geq\frac{(m-1)^2}{4}.
\]
Let \(\mathcal H\subset X\) be an isometrically embedded hyperbolic
horoball.  For every \(R>0\), the horoball contains a relatively compact
domain isometric to the radius-\(R\) ball
\(B_R^{\mathbb H}\subset\mathbb H^m(-1)\).  Transporting compactly
supported test functions from this domain to \(X\) gives
\[
  \lambda_1(X)
  \leq\lambda_1^D(B_R^{\mathbb H}).
\]
The balls \(B_R^{\mathbb H}\) exhaust hyperbolic space, and hence
\[
  \lambda_1^D(B_R^{\mathbb H})
  \longrightarrow
  \lambda_1\bigl(\mathbb H^m(-1)\bigr)
  =\frac{(m-1)^2}{4}.
\]
This proves the reverse inequality and therefore~\eqref{eq:cusp-spectrum}.
\end{proof}

\begin{remark}
\label{rem:cusp-p-spectrum}
The same argument applies for every \(1<p<\infty\).  Replacing McKean's
inequality by its \(p\)-version and using the variational
\(p\)-spectrum of hyperbolic space gives
\[
  \lambda_{1,p}(X)
  =
  \left(\frac{m-1}{p}\right)^p.
\]
Thus this cusp construction is not specific to the linear Laplacian.  The
corresponding defect identity and rigidity theorem are developed in
Appendix~\ref{app:p-rigidity}.
\end{remark}

Proposition~\ref{prop:hyperbolic-horoball-criterion} is only a criterion:
it neither produces a finite-volume quotient nor asserts that \(X\) is
nonhyperbolic.  The following construction provides both missing features
and has a two-sided sectional-curvature bound.

\begin{example}
\label{ex:bounded-curvature-finite-volume}
We construct the promised manifold.  It is complete, noncompact, and of
finite volume, with bounded negative sectional curvature; its universal
cover attains the sharp McKean value but is not hyperbolic.

\smallskip
\noindent\emph{Construction.}
Let \((M_0^m,g_0)\) be a complete noncompact finite-volume hyperbolic
manifold of constant sectional curvature \(-4\).  Choose one cusp and a
horospherical coordinate in which, for \(t\geq T\),
\[
  g_0=dt^2+\ee^{-4t}g_N,
\]
where \((N^{m-1},g_N)\) is compact and flat.  Choose a smooth
nonincreasing function
\[
  q\colon[T,\infty)\longrightarrow[1,2]
\]
that equals \(2\) near \(T\) and equals \(1\) for all sufficiently large
\(t\).  Set
\[
  f(t)
  =\ee^{-2T}
   \exp\left(-\int_T^tq(s)\,ds\right).
\]
Since \(q=2\) near \(T\), one has \(f(t)=\ee^{-2t}\) there.  Hence
\[
  g=dt^2+f(t)^2g_N
\]
fits smoothly with \(g_0\).  Replace \(g_0\) by this metric on the chosen
cusp and leave it unchanged elsewhere.  Denote the resulting manifold by
\((M,g)\).

\smallskip
\noindent\emph{Curvature bounds.}
Since \(g_N\) is flat and
\[
  \frac{f'}{f}=-q,
  \quad
  \frac{f''}{f}=q^2-q',
\]
the warped-product curvature formulas give, for a unit vector \(Y\)
tangent to \(N\) and a two-plane \(\sigma\) tangent to \(N\),
\[
  \sec_g(\partial_t,Y)=q'-q^2,
  \quad
  \sec_g(\sigma)=-q^2.
\]
Because \(1\leq q\leq2\) and \(q'\leq0\), both curvatures are at most
\(-1\).  Moreover, \(q'\) is bounded and
\[
  q'-q^2\geq-\|q'\|_\infty-4,
  \quad
  -q^2\geq-4,
\]
while outside the transition region the curvature is either \(-4\) or
\(-1\).  Therefore, with \(K^2=4+\|q'\|_\infty\),
\[
  -K^2\leq\sec_g\leq-1
  \quad\text{on \(M\)}.
\]

\smallskip
\noindent\emph{Finite volume and sharp spectrum.}
For all sufficiently large \(t\), one has \(q(t)=1\), and hence
\[
  f(t)=C\ee^{-t}
\]
for some \(C>0\).  The inclusion of the cusp cross-section
\(N\hookrightarrow M\) is \(\pi_1\)-injective.  Indeed, before the metric
change this is a cross-section of a hyperbolic cusp of \(M_0\), and
\(\pi_1(N)\) is its parabolic cusp subgroup; changing the metric does not
alter this topological inclusion.  Consequently, every component of the
inverse image of the far end in \(X=\widetilde M\) is
\[
  [T_1,\infty)\times\widetilde N
\]
with the lifted metric
\[
  dt^2+C^2\ee^{-2t}\widetilde g_N
\]
for some sufficiently large \(T_1\).  Since \(N\) is flat,
\((\widetilde N,\widetilde g_N)\) is Euclidean; after a linear change of
the horospherical coordinates, the displayed metric is the standard
curvature-\(-1\) metric on a hyperbolic horoball.  Also,
\[
  \vol_g\bigl([T,\infty)\times N\bigr)
  =\vol_{g_N}(N)\int_T^\infty f(t)^{m-1}\,dt
  <\infty.
\]
The metric is complete because the \(t\)-direction has infinite length.
Consequently, \((M,g)\) is complete and has finite volume.

Proposition~\ref{prop:hyperbolic-horoball-criterion} gives
\[
  \lambda_1(X)=\frac{(m-1)^2}{4}.
\]
On the other hand, \(g=g_0\) on the compact core, where the sectional
curvature is \(-4\).  Hence
\[
  X\not\cong\mathbb H^m(-1).
\]
This is the asserted bounded-curvature finite-volume failure of rigidity.
\end{example}

The quotient in Example~\ref{ex:bounded-curvature-finite-volume} has
uniformly bounded sectional curvature but not bounded geometry in the
standard sense.  Indeed, the lengths of nontrivial loops in the
horospherical cross-sections of the exact cusp decay exponentially, so
\(\operatorname{inj}_g\to0\) along the cusp.  Its universal cover \(X\),
on the other hand, has infinite injectivity radius and uniformly bounded
sectional curvature.

The minimizing functions used in
Proposition~\ref{prop:hyperbolic-horoball-criterion} may be supported in
balls moving arbitrarily far into the lifted cusp.  Their projections to
\(M\), and hence the associated probability measures on
\(Z\cong SM\), escape every compact subset.  Thus the defect identities
from Section~\ref{sec:busemann-defects} still hold, but the compactness step
in Proposition~\ref{prop:defect-support} no longer produces a limiting
probability measure.  This is exactly the compactness obstruction
identified at the beginning of the subsection.

Together, the warped-product construction and
Example~\ref{ex:bounded-curvature-finite-volume} identify the two losses
caused by removing closedness.  The first shows that an upper curvature
bound alone does not force the reference horosphere to be flat.  The
second shows that even finite volume and a two-sided curvature bound do
not prevent minimizing mass from escaping.  These examples do not exclude
rigidity under additional noncompact hypotheses: a suitable tightness
condition could replace compactness of the suspension, provided that the
lower curvature control needed in the final warped-product argument
continues to hold.

\begin{problem}
\label{prob:finite-volume-entropy-rigidity}
Let \((M^m,g)\) be a complete noncompact Riemannian manifold of finite
volume, and let \(X=\widetilde M\) be its universal Riemannian cover.
Suppose that
\[
  \Ric_X\geq-(m-1)g_X
\]
and
\[
  h_{\mathrm{vol}}(X)=m-1.
\]
Is \(X\) necessarily isometric to \(\mathbb H^m(-1)\)?
\end{problem}

This asks whether the volume-entropy rigidity theorem of
Ledrappier--X.~Wang, for which Gang Liu gave a short proof in the compact
case, extends to complete finite-volume quotients (see
\cite{LedrappierWang2010}, Theorem~2, and
\cite{Liu2011}, Theorem~1).  With bounded curvature and the additional
identity
\[
  \delta_{\pi_1(M)}=h_{\mathrm{vol}}(X)=m-1,
\]
the theorem of Barthelm\'e--Marquis--Zimmer gives
\(X\cong\mathbb H^m(-1)\) (see
\cite{BarthelmeMarquisZimmer2018}, Theorem~A).  Here
\(\delta_{\pi_1(M)}\) is the Poincar\'e exponent.  Its equality with
volume entropy is automatic for cocompact actions but can fail for
finite-volume nonuniform actions, even in pinched negative curvature
(see \cite{DalboPeignePicaudSambusetti2009}, Theorem~1.1).  To the best of our
knowledge, the problem remains open.  It is distinct from the spectral
equality studied here: an exact hyperbolic cusp can realize McKean's
spectral value without forcing the maximal volume-growth condition in the
problem.


\input{Section_7_Candel_detailed_20260728}

\appendix

\section{The \texorpdfstring{\(p\)}{p}-Laplacian rigidity theorem}
\label{app:p-rigidity}

This appendix proves Theorem~\ref{thm:p-rigidity}.  We treated the case
\(p=2\) in the main text because the quadratic defect displays the
geometry of the argument most clearly.  For general \(p\), its role is
played by the Bregman divergence associated with \(y\mapsto |y|^p\).

In the definition of \(\lambda_{1,p}(X)\), it suffices to use
nonnegative test functions.  Indeed, for \(v\in C_c^\infty(X)\), the
functions
\[
  F_\delta(v)=\sqrt{v^2+\delta^2}-\delta
\]
are smooth, nonnegative, and supported in \(\supp v\).  Moreover, they
converge to \(|v|\) in \(L^p\) and in \(p\)-energy as
\(\delta\downarrow0\), by dominated convergence and the bounds
\[
  |F_\delta(v)|\leq |v|,
  \qquad |\nabla F_\delta(v)|\leq |\nabla v|.
\]

An exact \(p\)-defect identity gives the lower bound, while equality
forces the defect to vanish.  For \(\rho=u^p\), this yields the same
linear transport equation as in the quadratic case, with coefficient
\(m-1\).  The leafwise-diffusion and support-saturation arguments
therefore apply without invoking a nonlinear \(p\)-heat flow.

\subsection{The nonlinear defect and critical sequences}

Fix \(p\in(1,\infty)\).  For two vectors \(x,y\) in the same
finite-dimensional inner-product space, define the convexity defect
\begin{equation}\label{eq:p-convexity-defect}
  \mc G_p(x,y)
  :=|x|^p+(p-1)|y|^p
    -p\langle |y|^{p-2}y,x\rangle,
\end{equation}
where \(|0|^{p-2}0\) is understood to be zero.
For \(p=2\), this reduces to
\(\mc G_2(x,y)=|x-y|^2\), the quadratic defect used in the main
argument.

\begin{lemma}\label{lem:p-uniform-convexity}
There is a constant \(c_p>0\), depending only on \(p\), such that
\(\mc G_p(x,y)\geq0\) for all \(x,y\), with equality if and only if
\(x=y\).  Moreover,
\[
  \mc G_p(x,y)\geq
  \begin{cases}
    c_p|x-y|^p,
      &p\geq2,\\[2mm]
    c_p(|x|+|y|)^{p-2}|x-y|^2,
      &1<p<2.
  \end{cases}
\]
The expression in the second line is defined to be zero when \(x=y=0\).
One may take
\[
  c_p=
  \begin{cases}
    p/4^p, & p\geq2,\\[1mm]
    p(p-1)/2, & 1<p<2.
  \end{cases}
\]
\end{lemma}

\begin{proof}
Let \(\Phi(z)=|z|^p\).  Then
\[
  \mc G_p(x,y)
  =\Phi(x)-\Phi(y)-D\Phi(y)[x-y],
\]
so the strict convexity of \(\Phi\) gives
\(\mc G_p(x,y)\geq0\), with equality precisely when \(x=y\).

Set \(e=x-y\).  There is nothing more to prove when \(e=0\), so assume
that \(e\neq0\).  Taylor's formula with integral remainder gives
\begin{equation}\label{eq:p-convexity-remainder}
  \mc G_p(x,y)
  =\int_0^1(1-t)
    D^2\Phi(y+te)[e,e]\,dt.
\end{equation}
For \(z\neq0\), decompose a vector \(v\) as
\[
  v=v_\perp+v_\parallel,
  \qquad
  v_\parallel=\frac{\langle v,z\rangle}{|z|^2}z,
  \qquad
  v_\perp\perp z.
\]
A direct computation gives
\[
  D^2\Phi(z)[v,v]
  =p|z|^{p-2}
    \bigl(|v_\perp|^2+(p-1)|v_\parallel|^2\bigr).
\]
If \(1<p<2\) and \(y+te=0\) at \(t=t_0\), the integrand in
\eqref{eq:p-convexity-remainder} grows at most like
\(|t-t_0|^{p-2}\).  Since \(p>1\), this singularity is integrable, and
we interpret \eqref{eq:p-convexity-remainder} as an improper integral.

\begin{enumerate}[label=\textup{(\roman*)}]
\item Suppose that \(p\geq2\).  Let
\(I=\{t\in[0,1]:|y+te|<|e|/4\}\).  This is an interval, and for
\(s,t\in I\),
\[
  |t-s||e|
  \leq |y+te|+|y+se|
  <\frac{|e|}{2},
\]
so \(I\) has length at most \(1/2\).  Hence
\(J=[0,3/4]\setminus I\) has measure at least \(1/4\), and every
\(t\in J\) satisfies \(1-t\geq1/4\) and
\(|y+te|\geq|e|/4\).
Since \(p-1\geq1\), the Hessian formula gives
\[
  D^2\Phi(z)[v,v]\geq p|z|^{p-2}|v|^2.
\]
Therefore \eqref{eq:p-convexity-remainder} yields
\[
  \mc G_p(x,y)
  \geq \int_J(1-t)p|y+te|^{p-2}|e|^2\,dt
  \geq \frac{p}{4^p}|e|^p.
\]

\item Suppose that \(1<p<2\).  Since \(p-1<1\), the Hessian formula
gives
\[
  D^2\Phi(z)[v,v]\geq p(p-1)|z|^{p-2}|v|^2.
\]
Moreover,
\[
  |y+te|
  =|(1-t)y+tx|
  \leq(1-t)|y|+t|x|
  \leq|x|+|y|.
\]
Since \(p-2<0\), it follows that
\[
  D^2\Phi(y+te)[e,e]
  \geq p(p-1)(|x|+|y|)^{p-2}|e|^2
\]
whenever \(y+te\neq0\).  Using the improper-integral interpretation at
a possible crossing, \eqref{eq:p-convexity-remainder} gives
\[
  \mc G_p(x,y)
  \geq\frac{p(p-1)}{2}(|x|+|y|)^{p-2}|e|^2,
\]
which proves the second estimate.
\end{enumerate}
\end{proof}

At the critical McKean scale, set
\begin{equation}\label{eq:p-critical-parameter}
  \alpha=\frac{m-1}{p}.
\end{equation}
Then \(p\alpha=m-1\), so the transport equation for the density
\(\rho=u^p\) has coefficient \(m-1\), independent of \(p\).

\begin{lemma}\label{lem:p-mckean-defect}
Let \((X^m,g)\) be a Riemannian manifold without boundary.  If
\(B\in C^2(X)\) satisfies \(|\nabla B|=1\), then every nonnegative
\(u\in C_c^\infty(X)\) satisfies the exact identity
\begin{equation}\label{eq:p-mckean-defect}
\begin{aligned}
  \int_X\bigl(|\nabla u|^p-\alpha^pu^p\bigr)\,dV
  ={}&\int_X\mc G_p(\nabla u,-\alpha u\nabla B)\,dV\\
     &+\alpha^{p-1}\int_X
       \bigl(\Delta B-(m-1)\bigr)u^p\,dV.
\end{aligned}
\end{equation}
\end{lemma}

\begin{proof}
Because \(u\geq0\), the function \(u^p\) belongs to \(C_c^1(X)\) and
\(\nabla(u^p)=pu^{p-1}\nabla u\).  Using \(|\nabla B|=1\), we obtain
the pointwise identity
\[
  \mc G_p(\nabla u,-\alpha u\nabla B)
  =|\nabla u|^p+(p-1)\alpha^pu^p
   +\alpha^{p-1}\langle\nabla B,\nabla(u^p)\rangle.
\]
Since \(u\) is compactly supported, integration by parts gives
\[
  \int_X\langle\nabla B,\nabla(u^p)\rangle\,dV
  =-\int_X(\Delta B)u^p\,dV.
\]
Consequently,
\[
\begin{aligned}
 &\int_X\mc G_p(\nabla u,-\alpha u\nabla B)\,dV
  +\alpha^{p-1}\int_X
    \bigl(\Delta B-(m-1)\bigr)u^p\,dV\\
 &\qquad
  =\int_X|\nabla u|^p\,dV
   +\bigl((p-1)\alpha^p-\alpha^{p-1}(m-1)\bigr)
      \int_Xu^p\,dV\\
 &\qquad
  =\int_X\bigl(|\nabla u|^p-\alpha^pu^p\bigr)\,dV,
\end{aligned}
\]
where the last equality uses \(m-1=p\alpha\).
\end{proof}

Assume now that \(\Delta B\geq m-1\).  Then
\eqref{eq:p-mckean-defect} expresses the spectral excess as the sum of
two nonnegative defects.  At the critical value
\(\lambda_{1,p}(X)=\alpha^p\), both defects must vanish along every
normalized minimizing sequence.  Lemma~\ref{lem:p-uniform-convexity}
turns the first vanishing defect into strong \(L^p\) control, and the
change of variables \(\rho=u^p\) converts this control into the linear
\(L^1\) transport equation used on the suspension.

\begin{proposition}\label{prop:p-critical-sequence}
Let \((X^m,g)\), \(m\geq2\), be complete, and let \(B\in C^2(X)\)
satisfy
\[
  |\nabla B|=1,
  \quad
  \Delta B\geq m-1,
\]
and suppose that
\[
  \lambda_{1,p}(X)=\alpha^p.
\]
Then there exists a nonnegative normalized minimizing sequence
\(\{u_j\}\subset C_c^\infty(X)\) satisfying
\[
  \int_Xu_j^p\,dV=1,
  \quad
  \int_X|\nabla u_j|^p\,dV\longrightarrow\alpha^p.
\]
Every such sequence has the following properties.
\begin{enumerate}[label=\textup{(\roman*)}]
\item Both defects vanish:
\begin{equation}\label{eq:p-two-defects}
\begin{split}
  \int_X\mc G_p(\nabla u_j,-\alpha u_j\nabla B)\,dV
  &\longrightarrow0,\\
  \int_X\bigl(\Delta B-(m-1)\bigr)u_j^p\,dV
  &\longrightarrow0.
\end{split}
\end{equation}
\item The first-order defect vanishes strongly in \(L^p(X)\):
\begin{equation}\label{eq:p-first-order-convergence}
  \|\nabla u_j+\alpha u_j\nabla B\|_{L^p(X)}
  \longrightarrow0.
\end{equation}
\item The probability densities \(\rho_j=u_j^p\) satisfy the
asymptotic transport equation
\begin{equation}\label{eq:p-density-transport}
  \|\nabla\rho_j+(m-1)\rho_j\nabla B\|_{L^1(X)}
  \longrightarrow0.
\end{equation}
\end{enumerate}
\end{proposition}

\begin{proof}
Existence follows from the reduction to nonnegative test functions at
the beginning of the appendix.  Fix any such sequence.  Applying
Lemma~\ref{lem:p-mckean-defect} and using the normalization, we obtain
\[
\begin{aligned}
  \int_X|\nabla u_j|^p\,dV-\alpha^p
  ={}&\int_X\mc G_p(\nabla u_j,-\alpha u_j\nabla B)\,dV\\
     &+\alpha^{p-1}\int_X
       \bigl(\Delta B-(m-1)\bigr)u_j^p\,dV.
\end{aligned}
\]
The left-hand side tends to zero, while both terms on the right are
nonnegative.  Since \(m\geq2\), we have \(\alpha>0\); hence both defect
integrals tend to zero.  This proves \textup{(i)}.

To prove \textup{(ii)}, set
\[
  R_j=\nabla u_j+\alpha u_j\nabla B,
  \quad
  A_j=|\nabla u_j|+\alpha u_j.
\]
\begin{enumerate}[label=\textup{(\arabic*)}]
\item If \(p\geq2\), the first estimate in
Lemma~\ref{lem:p-uniform-convexity} gives
\[
  c_p\|R_j\|_{L^p(X)}^p
  \leq
  \int_X\mc G_p(\nabla u_j,-\alpha u_j\nabla B)\,dV
  \longrightarrow0.
\]

\item If \(1<p<2\), the second estimate in
Lemma~\ref{lem:p-uniform-convexity} and \textup{(i)} give
\[
  I_j:=\int_XA_j^{p-2}|R_j|^2\,dV,
  \qquad
  c_pI_j
  \leq\int_X\mc G_p(\nabla u_j,-\alpha u_j\nabla B)\,dV
  \longrightarrow0.
\]
The integrand is interpreted as zero where \(A_j=0\), since then
\(R_j=0\).  Moreover,
\[
  \int_XA_j^p\,dV
  \leq 2^{p-1}
    \left(\int_X|\nabla u_j|^p\,dV+\alpha^p\right),
\]
so the sequence \(\{A_j\}\) is bounded in \(L^p(X)\).  H\"older's
inequality, with exponents \(2/p\) and \(2/(2-p)\), now gives
\[
\begin{aligned}
  \int_X|R_j|^p\,dV
  &=\int_X
    \bigl(A_j^{p-2}|R_j|^2\bigr)^{p/2}
    A_j^{p(2-p)/2}\,dV\\
  &\leq I_j^{p/2}
    \left(\int_XA_j^p\,dV\right)^{(2-p)/2}
  \longrightarrow0.
\end{aligned}
\]
\end{enumerate}
Thus \(R_j\to0\) in \(L^p(X)\) in both cases, proving \textup{(ii)}.

To prove \textup{(iii)}, set \(\rho_j=u_j^p\).  Since \(p>1\), the map
\(s\mapsto s^p\) is \(C^1\) on \([0,\infty)\), so
\(\rho_j\in C_c^1(X)\).  The chain rule and \(p\alpha=m-1\) give
\[
  \nabla\rho_j+(m-1)\rho_j\nabla B
  =pu_j^{p-1}R_j.
\]
Hence, by H\"older's inequality and \(\|u_j\|_{L^p}=1\),
\[
  \|\nabla\rho_j+(m-1)\rho_j\nabla B\|_{L^1(X)}
  \leq p\|u_j\|_{L^p(X)}^{p-1}\|R_j\|_{L^p(X)}
  =p\|R_j\|_{L^p(X)}
  \longrightarrow0.
\]
The last convergence follows from \textup{(ii)} and proves
\textup{(iii)}.
\end{proof}

\subsection{The sharp bound and the equality case}

\begin{proof}[Proof of Theorem~\ref{thm:p-rigidity}]
\emph{The lower bound.}
Fix \(p\in(1,\infty)\) and recall that \(\alpha=(m-1)/p\).
Since \(M\) is compact, its sectional curvature is also bounded below.
Thus there is \(a\geq1\) such that
\[
  -a^2\leq\sec_{\wti g}\leq-1.
\]
Proposition~\ref{prop:busemann-comparison} therefore gives, for every
\(\xi\in\partial_\infty\wti M\),
\[
  |\nabla B_\xi|=1,
  \quad
  \Delta B_\xi\geq m-1.
\]
Fix one such \(\xi\).  For every nonnegative
\(u\in C_c^\infty(\wti M)\), Lemma~\ref{lem:p-mckean-defect} and the
nonnegativity of both defects give
\[
  \int_{\wti M}|\nabla u|^p\,dV
  \geq\alpha^p\int_{\wti M}u^p\,dV.
\]
Taking the infimum and using the reduction to nonnegative test
functions yields
\[
  \lambda_{1,p}(\wti M)\geq\alpha^p.
\]

\emph{Rigidity in the equality case.}
Suppose that \(\lambda_{1,p}(\wti M)=\alpha^p\).
Fix \(\xi_0\in\partial_\infty\wti M\) and write
\(B_0=B_{\xi_0}\).  By
Proposition~\ref{prop:p-critical-sequence}, there is a nonnegative
normalized minimizing sequence \(u_j\).  Properties \textup{(i)} and
\textup{(iii)} of that proposition, with \(\rho_j=u_j^p\), give
\[
  \int_{\wti M}
    \bigl(\Delta B_0-(m-1)\bigr)\rho_j\,dV\longrightarrow0,
  \quad
  \|\nabla\rho_j+(m-1)\rho_j\nabla B_0\|_{L^1(\wti M)}
  \longrightarrow0.
\]
Let
\[
  Z=(\wti M\times\partial_\infty\wti M)/\Gamma,
  \quad
  \iota_{\xi_0}(x)=[x,\xi_0],
\]
and define probability measures on \(Z\) by
\[
  \mu_j=(\iota_{\xi_0})_*(\rho_jdV).
\]
Since \(Z\) is compact, a subsequence converges weakly to a probability
measure \(\mu\).  Because the Busemann defect \(D\) is continuous,
\[
  \int_ZD\,d\mu
  =\lim_{j\to\infty}\int_ZD\,d\mu_j
  =\lim_{j\to\infty}\int_{\wti M}
    \bigl(\Delta B_0-(m-1)\bigr)\rho_j\,dV
  =0.
\]
Since \(D\geq0\), it follows that
\begin{equation}\label{eq:p-limit-support}
  \supp\mu\subseteq D^{-1}(0).
\end{equation}

We next show that \(\supp\mu\) is leaf-saturated.  The transport
coefficient is \(m-1\), independent of \(p\), so the relevant operator
is the same linear leafwise operator used in
Section~\ref{sec:stationarity}:
\[
  L=\Delta^{\mathrm{leaf}}
    -(m-1)\langle V,\nabla^{\mathrm{leaf}}\,\cdot\,\rangle.
\]
Thus no nonlinear \(p\)-heat operator is needed.  For
\(f\in\mc D_0(A)\),
set \(f_0=f\circ\iota_{\xi_0}\).  Since
\(\rho_j\in C_c^1(\wti M)\) and \(f_0\in C^2(\wti M)\), integration by
parts for the compactly supported vector field
\(\rho_j\nabla f_0\) gives
\[
  \int_ZLf\,d\mu_j
  =-\int_{\wti M}
    \left\langle
      \nabla\rho_j+(m-1)\rho_j\nabla B_0,
      \nabla f_0
    \right\rangle dV.
\]
By Lemma~\ref{lem:classical-gradient-control},
\(\|\nabla f_0\|_\infty<\infty\).  Therefore
\[
  \left|\int_ZLf\,d\mu_j\right|
  \leq
  \|\nabla f_0\|_\infty
  \|\nabla\rho_j+(m-1)\rho_j\nabla B_0\|_{L^1(\wti M)}
  \longrightarrow0.
\]
Since \(Lf\) is continuous on \(Z\), weak convergence also gives
\[
  \int_ZLf\,d\mu=0
  \quad\text{for every }f\in\mc D_0(A).
\]
Lemma~\ref{lem:saturated-support}\textup{(ii)}--\textup{(iii)} now gives
stationarity of \(\mu\) and leaf saturation of its support.

Since \(\mu\) is a probability measure, its support is nonempty.  Choose
\(z_*=[x_*,\xi_*]\in\supp\mu\).  Leaf saturation implies that
\([x,\xi_*]\in\supp\mu\) for every \(x\in\wti M\).  Together with
\eqref{eq:p-limit-support}, this gives
\[
  [x,\xi_*]\in\supp\mu\subseteq D^{-1}(0).
\]
By the definition of \(D\), it follows that
\[
  \Delta B_{\xi_*}=m-1
\]
throughout \(\wti M\).  Moreover,
\(|\nabla B_{\xi_*}|=1\) by
Proposition~\ref{prop:busemann-comparison}.
Since \(M\) is closed, \(\sec_{\wti g}\) also has a uniform lower bound.
Lemma~\ref{lem:riemannian-rigidity} therefore gives
\[
  (\wti M,\wti g)\cong\bH^m(-1).
\]

\emph{The hyperbolic converse.}
Suppose that \(\wti M\cong\bH^m(-1)\).  The lower bound has already
been proved, so it remains only to establish the reverse inequality.
We use the radial-cutoff argument from the quadratic case, with
\(\alpha=(m-1)/p\).  Fix \(o\in\bH^m(-1)\), set \(r=d(o,\cdot)\), and
choose \(\eta\in C^\infty(\mathbb R)\) such that
\[
  0\leq\eta\leq1,\quad
  \eta(s)=0\ \text{for }s\leq0,\quad
  \eta(s)=1\ \text{for }s\geq1.
\]
For \(R>2\), define
\[
  \chi_R(r):=\eta(r-1)\eta(R+1-r),
  \quad
  u_R:=\chi_R(r)\ee^{-\alpha r}.
\]
Because \(\chi_R\) vanishes for \(r\leq1\), the nonsmoothness of the
distance function at \(o\) causes no problem, and
\(u_R\in C_c^\infty(\bH^m(-1))\).  Moreover,
\(\chi_R=1\) on \(2\leq r\leq R\), while its two transition regions
\([1,2]\) and \([R,R+1]\) have fixed width and
\(\|\chi_R'\|_\infty\) is independent of \(R\).

In polar coordinates,
\[
  dV=\sinh^{m-1}(r)\,dr\,d\omega,
\]
and \(p\alpha=m-1\) gives
\[
  \ee^{-p\alpha r}\sinh^{m-1}(r)
  =
  2^{-(m-1)}(1-\ee^{-2r})^{m-1}
  =
  2^{-(m-1)}+O(\ee^{-2r}).
\]
Hence, with
\(C_m=\vol(\mathbb S^{m-1})2^{-(m-1)}\),
\[
  \int_{\bH^m(-1)}u_R^p\,dV=C_mR+O(1).
\]
On \(2\leq r\leq R\), one has
\(|\nabla u_R|=\alpha u_R\).  On the transition regions,
\[
  |\nabla u_R|
  \leq(\|\chi_R'\|_\infty+\alpha)\ee^{-\alpha r},
\]
so their contribution to the \(p\)-energy is \(O(1)\), uniformly in
\(R\).  Therefore
\[
  \int_{\bH^m(-1)}|\nabla u_R|^p\,dV
  =\alpha^pC_mR+O(1).
\]
The Rayleigh quotients of \(u_R\) consequently tend to \(\alpha^p\), and
\[
  \lambda_{1,p}\bigl(\bH^m(-1)\bigr)\leq\alpha^p.
\]
Together with the lower bound, this proves
\[
  \lambda_{1,p}\bigl(\bH^m(-1)\bigr)
  =\left(\frac{m-1}{p}\right)^p
  \quad(1<p<\infty),
\]
and completes the proof.
\end{proof}

\begin{remark}\label{rem:p-spectrum-busemann}
A pointwise lower bound for \(\Delta B\) directly yields a lower bound
for the bottom of the \(p\)-spectrum.  More precisely, suppose that
\(X\) carries a function \(B\in C^2(X)\) such that \(|\nabla B|=1\)
and \(\Delta B\geq h>0\).  Applying
Lemma~\ref{lem:p-mckean-defect}, with \(m-1\) replaced by \(h\), gives
\[
  \lambda_{1,p}(X)\geq\left(\frac hp\right)^p.
\]
For the Busemann functions in Theorem~\ref{thm:p-rigidity},
Hessian comparison gives \(\Delta B_\xi\geq m-1\), and hence recovers
the McKean lower bound.
\end{remark}

\begin{bibdiv}
\begin{biblist}
\bibselect{references_20260727}
\end{biblist}
\end{bibdiv}

\end{document}
